\exercisesection{\S~4}

本段中所考察的李代数不一定是有限维的。

\begin{enumerate}
\def\labelenumi{\arabic{enumi})}
\tightlist
\item
  保留第2节的记号。令 \(\lambda \in k^{\mathrm{B}}\)。对任意 \(\alpha \in \mathrm{B}\)，将向量空间 E 的自同态 \(X_{-\alpha}^{\lambda}, H_{\alpha}^{\lambda}, X_{\alpha}^{\lambda}\) 与之关联，使得
\end{enumerate}

\[
X _ {- \alpha} ^ {\lambda} (\alpha_ {1}, \ldots , \alpha_ {n}) = (\alpha , \alpha_ {1}, \ldots , \alpha_ {n})
\]

\[
H _ {\alpha} ^ {\lambda} (\alpha_ {1}, \ldots , \alpha_ {n}) = (\lambda (\alpha) - \sum_ {i = 1} ^ {n} n (\alpha_ {i}, \alpha)) (\alpha_ {1}, \ldots , \alpha_ {n}).
\]

向量 \(X_{\alpha}^{\lambda}(\alpha_{1},\ldots,\alpha_{n})\) 由关于 n 的归纳用下式定义

\[
X _ {\alpha} ^ {\lambda} (\alpha_ {1}, \ldots , \alpha_ {n}) = (X _ {- \alpha_ {1}} ^ {\alpha} X _ {\alpha} ^ {\lambda} - \delta_ {\alpha , \alpha_ {1}} H _ {\alpha} ^ {\lambda}) (\alpha_ {2}, \ldots , \alpha_ {n}),
\]

其中 \(\delta_{\alpha, \alpha_1}\) 是 Kronecker 符号；约定，如果 \((\alpha_1, \ldots, \alpha_n)\) 是空字，则 \(X_{\alpha}^{\lambda}(\alpha_1, \ldots, \alpha_n)\) 为零。

证明引理1和2对这些自同态仍然成立。于是得到一个表示 \(\rho_{\lambda}:\mathfrak{a}\to \mathfrak{gl}(\mathrm{E})\)，使得

\[
\rho_ {\lambda} (x _ {\alpha}) = X _ {\alpha} ^ {\lambda}, \rho_ {\lambda} (h _ {\alpha}) = H _ {\alpha} ^ {\lambda}, \rho_ {\lambda} (x _ {- \alpha}) = X _ {- \alpha} ^ {\lambda}.
\]

¶ 2) 保留第3节的记号。设 m 是 a 的理想。

\begin{enumerate}
\def\labelenumi{\alph{enumi})}
\item
  令 \(\alpha\) 和 \(\beta\) 为 B 中两个不同的元素。假设当 N 充分大时，(ad \(x_{\alpha})^{\mathrm{N}}x_{\beta}\) 属于 \(\mathfrak{m}\)。证明 \(x_{\alpha \beta} \in \mathfrak{m}\)。（将§1第2节的结果应用于 \(\mathfrak{a} / \mathfrak{m}\)，其中赋予适当的 \(\mathfrak{sl}(2,k)\)-模结构。）
\item
  证明 \(n + \theta n\) 是 a 的最小有限余维理想。证明它也是包含 \((\operatorname{ad} x_{\alpha})^{4}x_{\beta}\) 和 \((\operatorname{ad} x_{-\alpha})^{4}x_{-\beta}\) 的最小理想。
\end{enumerate}

\begin{enumerate}
\def\labelenumi{\arabic{enumi})}
\setcounter{enumi}{2}
\tightlist
\item
  令 \((\mathfrak{g},\mathfrak{h})\) 为分裂半单李代数，R 为相应的根系，B 为 R 的一个基。对任意 \(\alpha \in \mathrm{B}\)（分别地，对任意一对 \((\alpha ,\beta)\in \mathrm{B}^2\)），令 \(\mathrm{R}(\alpha)\)（分别地，令 \(\mathrm{R}(\alpha ,\beta)\)）为 R 中由 \(\pm \alpha\) 构成的闭子集（分别地，为包含 \(\pm \alpha\) 和 \(\pm \beta\) 的最小闭子集）。令 \(\mathfrak{g}(\alpha)\)（分别地，令 \(\mathfrak{g}(\alpha ,\beta)\)）为代数 \(\mathfrak{h} + \mathfrak{g}^{\mathrm{R}(\alpha)}\)（分别地，代数 \(\mathfrak{h} + \mathfrak{g}^{\mathrm{R}(\alpha ,\beta)}\)）的导代数，参见§3。
\end{enumerate}

\begin{enumerate}
\def\labelenumi{\alph{enumi})}
\item
  证明 \(\mathfrak{g}(\alpha) = kH_{\alpha}\oplus \mathfrak{g}^{\alpha}\oplus \mathfrak{g}^{-\alpha}\)；它同构于 \(\mathfrak{sl}(2,k)\)。
\item
  证明 \(\mathfrak{g}(\alpha,\beta)\) 是半单的，且由 \(\mathfrak{g}(\alpha)\) 和 \(\mathfrak{g}(\beta)\) 生成。其根系可以与 \(\mathrm{R}(\alpha,\beta)\) 相识别。
\item
  令 \(\mathfrak{s}\) 为李代数（不一定有限维）。对所有 \(\alpha \in \mathrm{B}\)，令 \(f_{\alpha}\) 为从 \(\mathfrak{g}(\alpha)\) 到 \(\mathfrak{s}\) 的同态。假设对任意一对 \((\alpha, \beta)\)，存在同态 \(f_{\alpha \beta}: \mathfrak{g}(\alpha, \beta) \to \mathfrak{s}\)，同时延拓 \(f_{\alpha}\) 和 \(f_{\beta}\)。证明在这种情况下，存在唯一同态 \(f: \mathfrak{g} \to \mathfrak{s}\)，延拓各个 \(f_{\alpha}\)。（使用命题4 (i)。）
\end{enumerate}

\begin{enumerate}
\def\labelenumi{\arabic{enumi})}
\setcounter{enumi}{3}
\item
  令 \(\mathfrak{g}\) 为可分裂半单李代数，令 \(\sigma\) 为 \(k\) 的自同构。令 \(\mathfrak{g}_{\sigma}\) 为由 \(\mathfrak{g}\) 借助 \(\sigma\) 扩张标量所得到的李代数。证明 \(\mathfrak{g}_{\sigma}\) 同构于 \(\mathfrak{g}\)。（使用命题4的推论。）
\item
  \begin{enumerate}
  \def\labelenumii{\alph{enumii})}
  \tightlist
  \item
    令 g 为单李代数，令 \(k_{1}\) 为 g 的伴随表示的交换子。证明 \(k_{1}\) 是交换域，是 k 的有限扩张，并且 g 是绝对单的 \(k_{1}\)-李代数。
  \end{enumerate}
\end{enumerate}

反之，若 \(k_{1}\) 是 k 的有限扩张，且 g 是绝对单的 \(k_{1}\)-李代数，则 g 是单的 k-李代数，并且其伴随表示的交换子可以与 \(k_{1}\) 相识别。

\begin{enumerate}
\def\labelenumi{\alph{enumi})}
\setcounter{enumi}{1}
\tightlist
\item
  令 \(k'\) 为包含 \(k_{1}\) 的 Galois 扩张。证明 \(\mathfrak{g}_{(k')}\) 是 \([k_{1}:k]\) 个绝对单代数的直积。当 \(\mathfrak{g}_{(k')}\) 分裂时，这些代数彼此同构。（使用习题4。）
\end{enumerate}

\begin{enumerate}
\def\labelenumi{\arabic{enumi})}
\setcounter{enumi}{5}
\tightlist
\item
  令 \(\mathbf{A}\) 为交换环，令 \(\mathfrak{u}\) 为由生成族 \(\{x, y\}\) 及下列关系定义的 A-李代数
\end{enumerate}

\[
[ x, [ x, y ] ] = 0, \quad [ y, [ y, [ y, x ] ] ] = 0.
\]

证明 \(\mathfrak{u}\) 是以

\[
\{x, y, [ x, y ], [ y, [ x, y ] ] \}.
\]

证明当 \(\mathrm{A} = k\) 时，\(\mathfrak{u}\) 同构于 \(\mathfrak{a}_{+} / \mathfrak{n}\)，它与 \(\mathrm{B}_2\) 型根系相应。

¶7) 令 A 为其中 2 可逆的交换环，令 u 为由生成族 \(\{x,y\}\) 及下列关系定义的 A-李代数

\[
[ x, [ x, y ] ] = 0, \quad [ y, [ y, [ y, [ y, x ] ] ] ] = 0.
\]

证明 \(\mathfrak{u}\) 是以

\[
\left. \right.\left\{x, y, [ x, y ], [ y, [ x, y ] ], [ y, [ y, [ x, y ] ] ], [ x, [ y, [ y, [ x, y ] ] ] ] \right\}.
\]

证明当 A = k 时，u 同构于代数 \(a_{+}/n\)，它与 \(G_{2}\) 型根系相应。

\exercisesection{\S~5}

\begin{enumerate}
\def\labelenumi{\arabic{enumi})}
\item
  \(\operatorname{Aut}_{0}(\mathfrak{g})\) 在 \(\operatorname{Aut}(\mathfrak{g})\) 中的指数是有限的。
\item
  有 \(\mathrm{Aut}_e(\mathfrak{g}) = \mathrm{Aut}(\mathfrak{g})\)，当 \(\mathfrak{g}\) 可分裂、单且属于 \(\mathrm{G}_2, \mathrm{F}_4\) 或 \(\mathrm{E}_8\) 型时。
\item
  令 \(\mathfrak{h}\) 为 \(\mathfrak{g}\) 的分裂嘉当子代数，\(\mathfrak{b}\) 为 \((\mathfrak{g},\mathfrak{h})\) 的 Borel 子代数，\(\mathfrak{n} = [\mathfrak{b},\mathfrak{b}]\)，且 \(N = \exp \operatorname{ad}_{\mathfrak{g}}\mathfrak{n}\)。则
\end{enumerate}

\[
\operatorname{Aut} (\mathfrak {g}) = \mathrm{N}. \operatorname{Aut} (\mathfrak {g}, \mathfrak {h}). \mathrm{N}.
\]

(设 \(s \in \operatorname{Aut}(\mathfrak{g})\)。将§3第3节的命题10应用于 \(\mathfrak{b} \cap s(\mathfrak{b})\)，然后应用第VII章§3第4节的定理3。）

\begin{enumerate}
\def\labelenumi{\arabic{enumi})}
\setcounter{enumi}{3}
\tightlist
\item
  令 \(\mathfrak{h}\) 为 \(\mathfrak{g}\) 的嘉当子代数，令 \(s\) 为 \(\operatorname{Aut}(\mathfrak{g},\mathfrak{h})\) 中的一个元素，满足 \(sH \neq H\) 对所有非零 \(H\) 属于 \(\mathfrak{h}\) 成立。证明 \(s\) 的阶有限。（化归到 \(\mathfrak{h}\) 分裂的情形，并选取整数 \(n \geq 1\) 使得 \(\varepsilon(s)^n = 1\)。于是存在 \(\varphi \in \mathrm{T}_{\mathrm{Q}}\) 使得 \(f(\varphi) = s^n\)。令 \(\sigma\) 为 \(s|\mathfrak{h}\) 的转置。证明 \(1 + \sigma + \sigma^2 + \cdots + \sigma^{n-1} = 0\)，并由此推出 \(s^{n^2} = 1\)。）
\end{enumerate}

¶ 5) a) 令 \(a \in \operatorname{Aut}(\mathfrak{g})\)，令 \(\mathfrak{n}\) 为 \(a - 1\) 的幂零空间。证明下列条件等价：

\begin{enumerate}
\def\labelenumi{(\roman{enumi})}
\item
  \(\operatorname{Ker}(a - 1)\) 是 \(\mathfrak{g}\) 的嘉当子代数。
\item
  \(\mathfrak{n}\) 是 \(\mathfrak{g}\) 的嘉当子代数，且 \(a\in \mathrm{Aut}_0(\mathfrak{g})\)。
\item
  \(\dim\mathfrak{n}=\operatorname{rk}(\mathfrak{g})\)，且 \(a\in\operatorname{Aut}_{0}(\mathfrak{g})\)。
\end{enumerate}

\begin{enumerate}
\def\labelenumi{\alph{enumi})}
\setcounter{enumi}{1}
\item
  从现在起假设 k 是代数闭的。设 V 是向量空间，R 是 V 中的根系，\(T_{Q}\) 是群 \(\operatorname{Hom}(Q(R), k^{*})\)，n 是一个 \(\geq 1\) 的整数，而 \(T_{n}\) 是 \(T_{Q}\) 的子群，由阶整除 n 的元素组成。~设 \(\zeta\) 是本原 \(n\) 次单位根，属于 \(k\)。对于所有 \(H \in \mathrm{P}(\mathrm{R}^{\vee})\)，令 \(\psi(H)\) 为 \(\gamma \mapsto \zeta^{\gamma(H)}\) 所表示的 \(\mathrm{T}_{\mathrm{Q}}\) 中元素。映射 \(\psi\) 是从 \(\mathrm{P}(\mathrm{R}^{\vee})\) 到 \(\mathrm{T}_n\) 的同态，其核为 \(n\mathrm{P}(\mathrm{R}^{\vee})\)。令 \(t \in \mathrm{T}_n\)，令 \(\mathrm{C}\) 为 \(\mathrm{P}(\mathrm{R}^{\vee}) \otimes \mathbf{R}\) 中的一个单房。存在 \(w \in \mathrm{W}(\mathrm{R})\) 和 \(H \in \mathrm{P}(\mathrm{R}^{\vee})\)，使得 \(\frac{1}{n} H \in \overline{\mathrm{C}}\) 且 \(\psi(wH) = t\)。（使用第VI章§2第1节。）
\item
  令 \(\mathfrak{h}\) 为 \(\mathfrak{g}\) 的嘉当子代数，令 \(\mathbb{R}\) 为 \((\mathfrak{g},\mathfrak{h})\) 的根系。令 \(n,\zeta ,\psi\) 如 \(b\) ) 中，令 \(f\) 如第2节中。令 \(H\in \mathrm{P}(\mathrm{R}^{\vee})\)。在 \(\mathfrak{g}\) 中，\(f(\psi (H))\) 不变的元素集合为 \(\mathfrak{h}\oplus \sum_{\alpha \in \mathbb{R}'}\mathfrak{g}^\alpha\)，其中 \(\mathbb{R}'\) 是由 \(\alpha \in \mathbb{R}\) 且满足 \(\alpha (H)\in n\mathbf{Z}\) 的元素组成的集合；并且 \(f(\psi (H))\) 满足 \(a\) ) 中的条件，当且仅当 \(\frac{1}{n} H\) 属于一个单房。
\end{enumerate}

\(d\) ) 从现在起假设 \(\mathfrak{g}\) 是单的。令 \(\mathfrak{h}\) 和 \(\mathbb{R}\) 如 \(c\) ) 中，令 \((\alpha_{1},\ldots ,\alpha_{l})\) 为 \(\mathbb{R}\) 的一个基，令 \(h\) 为 \(\mathbb{R}\) 的 Coxeter 数，令 \(\zeta\) 为本原 \(h\) 次单位根，属于 \(k\)，并令 \(H\) 为 \(\mathfrak{h}\) 中满足 \(\alpha_{i}(H) = 1\)（\(i = 1,\dots,l\)）的元素。证明下列性质：

\begin{enumerate}
\def\labelenumi{(\roman{enumi})}
\item
  同态 \(\gamma \mapsto \zeta^{\gamma (H)}\) 从 Q(R) 到 \(k^{*}\) 定义了 \(\operatorname {Aut}(\mathfrak{g},\mathfrak{h})\) 中的一个元素，该元素满足 \(a\) ) 中的条件，且阶为 \(h\)。（使用 \(c\) )、第VI章§2的命题5和第VI章§1的命题31。）
\item
  \(\mathfrak{g}\) 的每个有限阶自同构都满足 \(a\) ) 中的条件，且其阶 \(\geq h\)。
\item
  \(\mathfrak{g}\) 的阶为 \(h\) 且满足 \(a\) ) 中条件的自同构在 \(\mathrm{Aut}_e(\mathfrak{g})\) 中构成一个共轭类。（使用第3节的命题5。）
\end{enumerate}

\begin{enumerate}
\def\labelenumi{\alph{enumi})}
\setcounter{enumi}{4}
\tightlist
\item
  令 \(\mathfrak{h}\) 和 \(\mathrm{R}\) 如 \(c\) ) 中，令 \(w\) 为 \(\mathrm{W(R)}\) 中的 Coxeter 变换。令 \(s \in \operatorname{Aut}(\mathfrak{g}, \mathfrak{h})\) 满足 \(\varepsilon(s) = w\)。证明 \(s\) 满足 \(a\) ) 中的条件，且阶为 \(h\)。（使用第VI章§1的命题33、第V章§6第2节和第VII章§4的命题9。）
\end{enumerate}

\(f)\) 若 \(s \in \operatorname{Aut}(\mathfrak{g})\)，则下列条件等价：

\begin{enumerate}
\def\labelenumi{(\roman{enumi})}
\item
  \(s\) 满足 \(a\) ) 中的条件，且阶为 \(h\)；
\item
  存在一个 \(\mathfrak{h}\)，它是 \(\mathfrak{g}\) 的在 \(s\) 下稳定的嘉当子代数，使得 \(s|\mathfrak{h}\) 是 \((\mathfrak{g},\mathfrak{h})\) 的 Weyl 群中的 Coxeter 变换。（使用 \(d\) ) 和 \(e\) )。）
\end{enumerate}

\(g\) ) \(d\) ) (i) 中的自同构的特征多项式为

\[
\mathrm{A} (\mathrm{T}) = (\mathrm{T} - 1) ^ {l} \prod_ {\alpha \in \mathrm{R}} (\mathrm{T} - \zeta^ {\alpha (H)}).
\]

自同构 \(s\) 在 \(e\) ) 中的特征多项式为

\[
\mathrm{B} (\mathrm{T}) = \left(\mathrm{T} ^ {h} - 1\right) ^ {l} \prod_ {i = 1} ^ {l} \left(\mathrm{T} - \zeta^ {m _ {i}}\right) \quad \left(m _ {i} \text { being   the   exponents   of } \mathrm{R}\right).
\]

（使用第VI章§1第11节的命题33 (iv)。）由关系 \(\mathrm{A(T)} = \mathrm{B(T)}\) 推出，对于所有 \(j \geq 1\)，满足该条件的 \(i\)（即 \(m_i \geq j\)）的个数等于满足条件 \(\alpha\in R_{+}\) 且 \(\alpha(H)=j\) 的元素个数；从而重新得到第VI章§4的习题6 c) 的结果。\(^{12}\)

\begin{enumerate}
\def\labelenumi{\arabic{enumi})}
\setcounter{enumi}{5}
\item
  假设 \(k = \mathbf{R}\) 或 \(\mathbf{C}\)。令 \(G\) 为李群 \(\mathrm{Aut}_0(\mathfrak{g})\)。证明元素 \(a \in G\) 是正则元（按第VII章§4第2节的意义）当且仅当它满足习题5 a) 的条件。
\item
  假设 g 是可分裂的。令 \(\mathrm{B}(\mathfrak{g})\) 为 g 的典则嘉当子代数的典则基（第3节，注2）。若 \(s \in \mathrm{Aut}(\mathfrak{g})\)，则 s 在 \(\mathrm{B}(\mathfrak{g})\) 上诱导一个置换；以 \(\operatorname{sgn}(s)\) 表示该置换的符号。证明 s 在 \(\bigwedge^{n} g\)（其中 \(n = \dim g\)）上的作用为
\end{enumerate}

\[
x \mapsto \operatorname{sgn} (s). x.
\]

¶8) 令 \(\bar{k}\) 为 k 的代数闭包，且 \(\bar{g} = \bar{k} \otimes_{k} g\)。Galois 群 \(\operatorname{Gal}(\bar{k}/k)\) 自然地作用于 \(\bar{g}\)、\(\bar{g}\) 的典则嘉当子代数（第3节，注2），以及其根系 \(\bar{R}\) 和典则基 \(\bar{B}\)。因而得到一个连续同态（即~具有开核）

\[
\pi : \operatorname{Gal} (\bar {k} / k) \to \operatorname{Aut} (\bar {\mathrm{R}}, \bar {\mathrm{B}}).
\]

证明 \(\mathfrak{g}\) 是可分裂的，当且仅当满足下列两个条件：

\begin{enumerate}
\def\labelenumi{(\roman{enumi})}
\item
  同态 \(\pi\) 是平凡的。
\item
  g 有一个 Borel 子代数 b。
\end{enumerate}

(证明，\(\mathfrak{h}\) 包含于 \(\mathfrak{b}\) 中的情形是分裂的，当且仅当 \(\pi\) 是平凡的。)

\begin{enumerate}
\def\labelenumi{\arabic{enumi})}
\setcounter{enumi}{8}
\tightlist
\item
  令 R 为既约根系，B 为 R 的一个基，且 \((\mathfrak{g}_{0}, \mathfrak{h}_{0}, \mathrm{B}, (X_{\alpha})_{\alpha \in \mathrm{B}})\) 为相应的带标架半单李代数（§4）。令 \(\bar{k}\) 为 k 的代数闭包，令 \(\rho : \operatorname{Gal}(\bar{k}/k) \to \operatorname{Aut}(\mathrm{R}, \mathrm{B})\) 为连续同态（参见~习题8）；若 \(\sigma \in \operatorname{Gal}(\bar{k}/k)\)，以 \(\rho_{\sigma}\) 表示 \(\bar{g}_{0} = \bar{k} \otimes_{k} g_{0}\) 的 k-线性自同构，使得 \(\rho_{\sigma}(X_{\alpha}) = X_{\rho(\sigma)\alpha}\)。另一方面，\(\operatorname{Gal}(\bar{k}/k)\) 在 \(\bar{k}\) 上的自然作用可以延拓为在 \(\bar{g}_{0}\) 上的作用。令 g 为 \(\bar{g}_{0}\) 的子集，由满足 \(\rho_{\sigma}(x) = \sigma^{-1}.x\) 对所有 \(\sigma \in \operatorname{Gal}(\bar{k}/k)\) 成立的元素 x 构成。
\end{enumerate}

\begin{enumerate}
\def\labelenumi{\alph{enumi})}
\item
  证明 \(\mathfrak{g}\) 是 \(k\)-李子代数，包含于 \(\bar{\mathfrak{g}}_0\)；并且 \(\mathfrak{g}\) 到 \(\bar{\mathfrak{g}}_0\) 的嵌入延拓为从 \(\bar{k} \otimes_k \mathfrak{g}\) 到 \(\bar{\mathfrak{g}}_0\) 的一个同构。特别地，\(\mathfrak{g}\) 是半单的。
\item
  令 \(\mathfrak{b}_0\) 为 \(\mathfrak{g}_0\) 的子代数，它由 \(\mathfrak{h}_0\) 和 \(X_{\alpha}\) 生成。置
\end{enumerate}

\[
\bar {\mathfrak {b}} _ {0} = \bar {k} \otimes_ {k} \mathfrak {b} _ {0}, \quad \mathfrak {h} = \mathfrak {g} \cap \bar {\mathfrak {b}} _ {0}, \quad \bar {\mathfrak {h}} _ {0} = \bar {k} \otimes_ {k} \mathfrak {h} _ {0}, \quad \mathfrak {h} = \mathfrak {g} \cap \bar {\mathfrak {h}} _ {0}.
\]

证明 \(\bar{k}\otimes_{k}b=\bar{b}_{0}\) 且 \(\bar{k}\otimes_{k}h=\bar{h}_{0}\)，从而 b 是 g 的 Borel 子代数，而 h 是包含于其中的嘉当子代数。

\begin{enumerate}
\def\labelenumi{\alph{enumi})}
\setcounter{enumi}{2}
\tightlist
\item
  证明与代数 g 相关联的同态 \(\pi\)（参见~习题8）等于 \(\rho\)。
\end{enumerate}

¶ 10) 令 h 为 g 的分裂嘉当子代数，且 \((X_{\alpha})_{\alpha\in\mathbb{R}}\) 为 \((\mathfrak{g},\mathfrak{h})\) 中的谢瓦莱系，参见~§2第4节。若 \(\alpha\in R\)，置

\[
\theta_ {\alpha} = e ^ {\mathrm{ad} X _ {\alpha}} e ^ {\mathrm{ad} X _ {- \alpha}} e ^ {\mathrm{ad} X _ {\alpha}},
\]

并以 \(\overline{W}\) 表示 \(\operatorname{Aut}_{e}(\mathfrak{g},\mathfrak{h})\) 中由 \(\theta_{\alpha}\) 生成的子群。

\begin{enumerate}
\def\labelenumi{\alph{enumi})}
\item
  证明 \(\varepsilon (\overline{\mathrm{W}}) = \mathrm{W}(\mathrm{R})\)。
\item
  令 \(s \in \overline{\mathrm{W}}\)，令 \(w = \varepsilon(s)\)。证明
\end{enumerate}

\[
s (X _ {\alpha}) = \pm X _ {w (\alpha)} \quad \text {   for   all   } \alpha \in \mathrm{R}.
\]

（使用§2的习题5。）

\begin{enumerate}
\def\labelenumi{\alph{enumi})}
\setcounter{enumi}{2}
\tightlist
\item
  令 M 为 \(\varepsilon : \overline{\mathrm{W}} \to \mathrm{W}(\mathrm{R})\) 的核。证明 M 包含于 \(f(\mathrm{T}_{\mathrm{Q}})\) 的子群，该子群由元素 \(f(\varphi)\) 组成，其中 \(\varphi^2 = 1\)。证明 M 包含由下式定义的元素 \(f(\varphi_{\alpha})\)，其中 \(\varphi_{\alpha}(\beta) = (-1)^{\langle \beta, \alpha^{\vee} \rangle}\)（注意 \(\theta_{\alpha}^{2} = f(\varphi_{\alpha})\)。）
\end{enumerate}

\(d)\) *令 \(\varphi\in\mathrm{Hom}(\mathbf{Q},\{\pm1\})\)。证明 \(f(\varphi)\) 属于 M，当且仅当 \(\varphi\) 延拓为一个从 P 到 \(\{\pm1\}\) 的同态。（充分性由 M 包含 \(f(\varphi_{\alpha})\) 这一事实得到。为证明必要性，化归到 \(k=\mathbf{Q}\) 的情形，并使用 M 包含于 \(f(\mathrm{T}_{\mathrm{Q}})\cap\mathrm{Aut}_{e}(\mathfrak{g})=\mathrm{Im}(\mathrm{T}_{\mathrm{P}})\) 的事实，参见~§7，习题26 \(d)\)。推出 M 同构于群 \(\mathrm{Q}/(\mathrm{Q}\cap2\mathrm{P}).^{13}\) 的对偶。

\begin{enumerate}
\def\labelenumi{\arabic{enumi})}
\setcounter{enumi}{10}
\tightlist
\item
  按第2节的记号，假设 \(k\) 是非离散且局部紧的，因而同构于 \(\mathbf{R}\)、\(\mathbf{C}\) 或 \(\mathbf{Q}_p\) 的有限扩张（《交换代数》，第VI章§9第3节）。对所有 \(n \geq 1\)，商 \(k^* / k^{*n}\) 是有限的（超度量情形参见~《交换代数》第VI章§9的习题3）。推出商 \(T_Q / \text{Im}(T_P)\) 和 \(\text{Aut}(\mathfrak{g}) / \text{Aut}_e(\mathfrak{g})\) 是有限的。
\end{enumerate}

当 \(k = \mathbf{R}\) 时，证明 \(\mathrm{T_Q / Im(T_P)}\) 同构于 \(\mathbf{F}_{2}\)-向量空间 \((\mathrm{Q} \cap 2\mathrm{P}) / 2\mathrm{Q}\) 的对偶。当 \(k = \mathbf{C}\) 时，有 \(\mathrm{T_Q = Im(T_P)}\)；这是李群 \(\operatorname{Aut}(\mathfrak{g})\) 中李代数为 \(\mathfrak{h}\) 的整子群。

¶ 12) 令 \(\mathfrak{h}\) 为 \(\mathfrak{g}\) 的分裂嘉当子代数，A 为 \(\mathfrak{h}\) 的子集，且 \(s \in \operatorname{Aut}(\mathfrak{g})\) 满足 \(s\mathrm{A} = \mathrm{A}\)。证明存在 \(t \in \operatorname{Aut}(\mathfrak{g}, \mathfrak{h})\)，使得 \(t|\mathrm{A} = s|\mathrm{A}\) 且 \(ts^{-1} \in \operatorname{Aut}_e(\mathfrak{g})\)。（令 \(\mathfrak{a}\) 为 A 在 \(\mathfrak{g}\) 中的交换子；这是 \(\mathfrak{g}\) 的约化子代数，其中 \(sh\) 和 \(\mathfrak{h}\) 是分裂嘉当子代数；推出存在 \(u \in \operatorname{Aut}_e(\mathfrak{a})\)，使得 \(ush = \mathfrak{h}\)；证明存在 \(v \in \operatorname{Aut}_e(\mathfrak{g})\) 延拓 \(u\)，使得 \(v|\mathrm{A} = \mathrm{Id}_{\mathrm{A}}\)；取 \(t = vs.\)）推出，如果 \(s \in \operatorname{Aut}_0(\mathfrak{g})\)，则存在 \(w \in \mathrm{W}(\mathrm{R})\)，使得 \(w|\mathrm{A} = s|\mathrm{A}\)。

¶13) 令 \((\mathfrak{g},\mathfrak{h},\mathrm{B},(X_{\alpha})_{\alpha\in\mathrm{B}})\) 为带标架半单李代数，R 为 \((\mathfrak{g},\mathfrak{h})\) 的根系，\(\Delta\) 为相应的邓金图，且 \(\varPhi\) 为下式的一个子群

\[
\operatorname{Aut} (\mathrm{R}, \mathrm{B}) = \operatorname{Aut} (\Delta).
\]

若 \(s \in \varPhi\)，则按以下条件将 \(s\) 延拓为 \(\mathfrak{g}\) 的一个自同构

\[
s (X _ {\alpha}) = X _ {s \alpha}, \quad \text { and } \quad s (H _ {\alpha}) = H _ {s \alpha} \quad \text { for   all } \alpha \in \mathrm{B}, \text { cf.   Prop. } 1.
\]

于是将 \(\varPhi\) 与 \(\operatorname{Aut}(\mathfrak{g},\mathfrak{h})\) 的一个子群相识别；以 \(\tilde{\mathfrak{g}}\)（分别地，\(\tilde{\mathfrak{h}}\)）表示 g（分别地，h）中在 \(\varPhi\) 下不变的元素所成的子代数。

\begin{enumerate}
\def\labelenumi{\alph{enumi})}
\tightlist
\item
  令 \(\alpha \in \mathbf{B}\)，令 \(\mathbf{X} = \varPhi.\alpha\)。利用第VI章中的图版，证明只有以下两种情形可能：
\end{enumerate}

\begin{enumerate}
\def\labelenumi{(\roman{enumi})}
\item
  \(X\) 中异于 \(\alpha\) 的每个元素都与 \(\alpha\) 正交；
\item
  存在唯一元素 \(\alpha'\)，它属于 \(X - \{\alpha\}\)，且不与 \(\alpha\) 正交，并且 \(n(\alpha, \alpha') = n(\alpha', \alpha) = -1\)。
\end{enumerate}

\begin{enumerate}
\def\labelenumi{\alph{enumi})}
\setcounter{enumi}{1}
\tightlist
\item
  令 i 为限制映射 \(h^{*} \to \tilde{h}^{*}\)，令 \(\tilde{\mathbf{B}} = i(\mathbf{B})\)；映射 \(\mathbf{B} \to \tilde{\mathbf{B}}\) 将 \(\tilde{B}\) 与 \(B/\varPhi\) 相识别。证明 \(\tilde{g}\) 是半单李代数，\(\tilde{h}\) 是其中的分裂嘉当子代数，且 \(\tilde{B}\) 是 \(\mathrm{R}(\tilde{\mathfrak{g}}, \tilde{\mathfrak{h}})\) 的一个基。（注意 \(\tilde{B}\) 包含于 \(\mathrm{R}(\tilde{\mathfrak{g}}, \tilde{\mathfrak{h}})\) 中，且 \(\mathrm{R}(\tilde{\mathfrak{g}}, \tilde{\mathfrak{h}})\) 的每个元素都是 \(\tilde{B}\) 中元素的具有相同符号的整数系数线性组合。）若 \(\tilde{\alpha} \in \tilde{B}\)，相应的逆根 \(H_{\tilde{\alpha}} \in \tilde{h}\) 由下式给出
\end{enumerate}

\[
\begin{array}{l} H _ {\tilde {\alpha}} = \sum_ {i (\alpha) = \tilde {\alpha}} H _ {\alpha} \quad \text {in case (i) of a)} \\ H _ {\tilde {\alpha}} = 2 \sum_ {i (\alpha) = \tilde {\alpha}} H _ {\alpha} \quad \text {in case (ii) of a)}, \end{array}
\]

其中求和遍历 \(\alpha \in \mathrm{B}\) 中满足 \(i(\alpha) = \tilde{\alpha}\) 的元素。若 \(\beta \in \mathrm{B}\) 的像为 \(\tilde{\beta} \in \tilde{\mathrm{B}}\)，则

\[
  n (\tilde {\beta}, \tilde {\alpha}) = \sum_ {i (\alpha) = \tilde {\alpha}} n (\beta , \alpha) \quad \text { in   case   (i) }
\]

\[
n (\tilde {\beta}, \tilde {\alpha}) = 2 \sum_ {i (\alpha) = \tilde {\alpha}} n (\beta , \alpha) \quad \text { in   case   (ii) }.
\]

推出 \(\mathrm{R}(\tilde{\mathfrak{g}},\tilde{\mathfrak{h}})\) 的邓金图如何由二元组 \((\varDelta,\varPhi)\) 确定。

\begin{enumerate}
\def\labelenumi{\alph{enumi})}
\setcounter{enumi}{2}
\tightlist
\item
  证明若 \(\mathfrak{g}\) 是单的，则 \(\tilde{\mathfrak{g}}\) 也是单的。
\end{enumerate}

若 \(\mathfrak{g}\) 属于 \(A_l\) 型，\(l \geq 2\)，且 \(\varPhi\) 的阶为2，则 \(\tilde{\mathfrak{g}}\) 属于 \(B_{l/2}\) 型（当 \(l\) 为偶数时），属于 \(C_{(l+1)/2}\) 型（当 \(l\) 为奇数时）。

若 \(\mathfrak{g}\) 属于 \(D_l\) 型，\(l \geq 4\)，且 \(\varPhi\) 的阶为2，则 \(\tilde{\mathfrak{g}}\) 属于 \(B_{l-1}\) 型。

若 \(\mathfrak{g}\) 属于 \(\mathrm{D}_4\) 型，且 \(\varPhi\) 的阶为3或6，则 \(\tilde{\mathfrak{g}}\) 属于 \(\mathrm{G}_2\) 型。

若 g 属于 \(E_{6}\) 型，且 \(\Phi\) 的阶为2，则 \(\tilde{g}\) 属于 \(F_{4}\) 型。

\exercisesection{\S~6}

\begin{enumerate}
\def\labelenumi{\arabic{enumi})}
\item
  证明 \(Z(\lambda)\) 可以定义为§4习题1中的表示 \(\rho_{\lambda}\) 的一个商。
\item
  令 \(\mu\) 为 \(Z(\lambda)\)（分别地，\(E(\lambda)\)）的一个权。证明存在权序列 \(\mu_0, \ldots, \mu_n\)，属于 \(Z(\lambda)\)（分别地，\(E(\lambda)\)），使得 \(\mu_0 = \lambda, \mu_n = \mu\)，且 \(\mu_{i-1} - \mu_i \in B\) 对 \(1 \leq i \leq n\) 成立。
\item
  假设 g 是单的，并以 \(\tilde{\alpha}\) 表示 R 的最高根。模 \(\mathrm{E}(\tilde{\alpha})\) 同构于赋有伴随表示的 g。若 C 是与 g 的 Killing 型相关联的 Casimir 元，则 C 在 \(\mathrm{End}\mathrm{E}(\tilde{\alpha})\) 中的像为单位元（参见~第I章§3第7节的命题12）。利用命题7的推论推出
\end{enumerate}

\[
\varPhi_ {\mathrm{R}} (\tilde {\alpha}, \tilde {\alpha} + 2 \rho) = 1,
\]

其中 \(\Phi_{\mathrm{R}}\) 是 \(\mathfrak{h}^*\) 上的典则双线性型（第VI章§1第12节）。

\begin{enumerate}
\def\labelenumi{\arabic{enumi})}
\setcounter{enumi}{3}
\tightlist
\item
  沿用第4节的记号。
\end{enumerate}

\begin{enumerate}
\def\labelenumi{\alph{enumi})}
\item
  令 \(m \in \mathbf{N}\)。给 \(\mathbf{N}^m\) 赋予乘积序。证明对于任意 \(S\)，若它是 \(\mathbf{N}^m\) 的子集，则 \(S\) 的极小元素集是有限的。
\item
  令 \(\alpha_{1},\ldots ,\alpha_{m}\) 为 R 中互异的元素，\(X_{i}\in \mathfrak{g}^{\alpha_{i}} - \{0\}\)，令 S 为非零序列 \((p_i)\in \mathbf{N}^m\) 且满足 \(\sum p_i\alpha_i = 0\) 的集合，M 为 S 的极小元素集。则 \(\mathfrak{h}\) 与 \(X_1^{p_1}\dots X_m^{p_m}\)（其中 \((p_1,\dots ,p_m)\in \mathrm{M}\)）生成代数 \(\mathrm{U}^0\)。
\item
  证明 \(U^{0}\) 既是左诺特代数又是右诺特代数。（给 \(U^{0}\) 赋予由 \(\mathrm{U}(\mathfrak{g})\) 的滤过诱导的滤过，并利用 a) 和 b) 证明 gr \(U^{0}\) 是有限型交换代数。）
\end{enumerate}

\(d\) ) 证明对于所有 \(\lambda \in \mathfrak{h}^*\)，\(\mathrm{U}^{\lambda}\) 是有限型的左（分别地，右）\(\mathrm{U}^0\)-模。

\begin{enumerate}
\def\labelenumi{\alph{enumi})}
\setcounter{enumi}{4}
\tightlist
\item
  令 V 为满足 \(V = \bigoplus_{\lambda \in b^{*}} V_{\lambda}\) 的单 g-模。若其中某个 \(V_{\lambda} \neq 0\) 是有限维的，则所有的 \(V_{\lambda}\) 都是有限维的。（使用 d）。
\end{enumerate}

\begin{enumerate}
\def\labelenumi{\arabic{enumi})}
\setcounter{enumi}{4}
\tightlist
  \item
  证明若 \(\mathfrak{g} = \mathfrak{sl}(2,k)\)，则本段的模 \(Z(\lambda)\) 同构于§1习题2中的模 \(Z(\lambda)\)。
\end{enumerate}

\exercisesection{\S~7}

所考察的所有 g-模（习题14和15中的模除外）均假定为有限维的。

\begin{enumerate}
\def\labelenumi{\arabic{enumi})}
\tightlist
\item
  令 \(\omega \in \mathrm{P}_{++}\)；以 \(\mathrm{S}(\omega)\) 表示 \(\mathrm{E}(\omega)\) 的权集，换言之，它是包含 \(\omega\) 的最小 R-饱和子集（命题5）。若 \(\lambda \in \mathrm{S}(\omega)\)，则有 \(\lambda \equiv \omega\)（mod. Q）。反之，令 \(\lambda \in \mathrm{P}\) 满足 \(\lambda \equiv \omega\)（mod. Q)；证明下列性质等价：
\end{enumerate}

\begin{enumerate}
\def\labelenumi{(\roman{enumi})}
\item
  \(\lambda \in S(\omega);\)
\item
  \(\omega - w\lambda \in Q_{+}\) 对所有 \(w \in W\) 均成立；
\item
  \(\lambda\) 属于 W 的 \(\omega\) 在 \(\mathfrak{h}_{\mathbf{R}}^{*}\) 中的凸包。
\end{enumerate}

（为证明 (iii) \(\Longrightarrow\) (ii)，注意 \(\omega - w\omega\) 是 \(\mathrm{R}_{+}\) 中元素的系数 \(\geq 0\) 的线性组合；由凸性推出 \(\omega - w\lambda\) 也有同样的性质；由于 \(\omega - w\lambda\) 属于 Q，这意味着 \(\omega - w\lambda \in \mathrm{Q}_{+}\)。为证明 (ii) \(\Longrightarrow\) (i)，选取 \(w\) 使得 \(w\lambda \in \mathrm{P}_{++}\)，并应用命题3的推论2。(i) \(\Longrightarrow\) (iii) 的蕴含是直接的。）

\begin{enumerate}
\def\labelenumi{\arabic{enumi})}
\setcounter{enumi}{1}
\tightlist
\item
  令 \((\mathrm{R}_{i})_{i\in\mathrm{I}}\) 为 R 的不可约分支族，且 \(g=\prod_{i\in I}g_{i}\) 为 g 分解为单代数之积的相应分解。将 P 与 \(\mathrm{P}(\mathrm{R}_{i})\) 的乘积相识别，并给每个 \(\mathrm{P}(\mathrm{R}_{i})\) 赋予由基 \(B_{i}=B\cap R_{i}\) 定义的序关系。
\end{enumerate}

\begin{enumerate}
\def\labelenumi{\alph{enumi})}
\item
  令 \(\omega = (\omega_{i})_{i\in \mathrm{I}}\) 为 \(\mathrm{P}_{++} = \prod_{i\in \mathrm{I}}\mathrm{P}_{++}(\mathrm{R}_i)\) 中的一个元素。证明单 \(\mathfrak{g}\)-模 \(\operatorname {E}(\omega)\) 同构于单 \(\mathfrak{g}_i\)-模 \(\operatorname {E}(\omega_i)\) 的张量积。
\item
  令 \(\mathcal{M}\)（分别地，\(M_{i}\)）为 \(P_{++}\)（分别地，\(P_{++}(R_{i})\)）中具有命题6和7的 (i)、(ii)、(iii)、(iv) 的等价性质的元素集。证明 \(M = \prod_{i \in I} M_{i}\)，换言之，\(\omega \in M\) 当且仅当对于所有 \(i \in I\)，\(\omega_{i}\) 或为零，或为 \(R_{i}\) 的极小权。推出 M 是 P/Q 中元素在 P 里的代表元系。
\item
  令 E 为单 g-模，X 为其权集。证明 X 含有 M 中唯一的一个元素，且该元素的重数等于 X 中各元素重数的上界。
\end{enumerate}

\begin{enumerate}
\def\labelenumi{\arabic{enumi})}
\setcounter{enumi}{2}
\tightlist
\item
  \begin{enumerate}
  \def\labelenumii{\alph{enumii})}
  \tightlist
  \item
    令 E 为 g-模。证明下列条件等价：
  \end{enumerate}
\end{enumerate}

\begin{enumerate}
\def\labelenumi{(\roman{enumi})}
\item
  \(\mathfrak{g}\) 与 E 的半直积的秩严格大于 \(\mathfrak{g}\) 的秩。
\item
  0 是 E 的一个权。
\item
  \(\mathrm{E}\) 存在一个根权（即~属于 \(\mathrm{Q}\) 的权）。
\end{enumerate}

\begin{enumerate}
\def\labelenumi{\alph{enumi})}
\setcounter{enumi}{1}
\tightlist
\item
  假设 \(\mathrm{E}\) 是单的。证明 (i)、(ii)、(iii) 等价于
\end{enumerate}

\begin{enumerate}
\def\labelenumi{(\roman{enumi})}
\setcounter{enumi}{3}
\tightlist
\item
  E 的最高权是根权。
\end{enumerate}

若这些条件满足，则不存在非零不变交错双线性型。（使用命题12和§6的命题1 (ii)。）

\begin{enumerate}
\def\labelenumi{\arabic{enumi})}
\setcounter{enumi}{3}
\item
  令 \(k'\) 为 \(k\) 的一个扩张，且 \(\mathfrak{g}' = \mathfrak{g}_{(k')}\)。证明每个 \(\mathfrak{g}'\)-模都由某个 \(\mathfrak{g}\)-模经扩张标量得到，且该模在同构意义下唯一。
\item
  令 \(\mathbf{E}\) 为 \(\mathfrak{g}\)-模。证明下列条件等价：
\end{enumerate}

\begin{enumerate}
\def\labelenumi{(\roman{enumi})}
\item
  E 是忠实的（即~从 g 到 gl(E) 的典则映射是单射）。
\item
  g 的每个根都是 E 的两个权之差。
\end{enumerate}

¶6) 令 \(\varphi\) 为 \(\mathfrak{g}\) 的一个对合自同构，其在 \(\mathfrak{h}\) 上的限制为 -Id。证明若 E 是 \(\mathfrak{g}\)-模，则 E 上存在非退化对称双线性型 \(\varPsi\)，使得

\[
\varPsi (x. a, b) + \varPsi (a, \varphi (x). b) = 0 \quad \text {   for   } x \in \mathfrak {g}, a, b \in \mathrm{E}.
\]

（化归到 E 为单的情形。证明 E 经 \(\varphi\) 变换所得的模同构于 E 的对偶 \(E^{*}\)，从而推出存在满足上述条件的非退化双线性型 \(\Psi\)。接着证明若 e 是 E 中的本原元，则 \(\Psi(e,e)\neq0\)。推出 \(\Psi\) 是对称的。）

\begin{enumerate}
\def\labelenumi{\arabic{enumi})}
\setcounter{enumi}{6}
\tightlist
\item
  若 \(\lambda \in \mathrm{P}_{++}\)，以 \(\rho_{\lambda}:\mathrm{U}(\mathfrak{g})\to \mathrm{End}(\mathrm{E}(\lambda))\) 表示由单模 \(\mathrm{E}(\lambda)\) 定义的表示。我们有 \(\mathrm{Im}(\rho_{\lambda}) = \mathrm{End}(\mathrm{E}(\lambda));\) 置 \(\mathfrak{m}_{\lambda} = \mathrm{Ker}(\rho_{\lambda})\)。\(a)\) 证明 \(\mathfrak{m}_{\lambda}\) 两两不同，并且它们正是这样的所有双边理想：\(\mathfrak{m}\) 属于 \(\mathrm{U}(\mathfrak{g})\)，且 \(\mathrm{U}(\mathfrak{g}) / \mathfrak{m}\) 为有限维单 \(k\)-代数。
\end{enumerate}

\begin{enumerate}
\def\labelenumi{\alph{enumi})}
\setcounter{enumi}{1}
\item
  若 I 是 \(P_{++}\) 的有限子集，置 \(m_{I} = \bigcap_{\lambda \in I} m_{\lambda}\)。证明典则映射 \(U(\mathfrak{g}) / \mathfrak{m}_{I} \to \prod_{\lambda \in I} U(\mathfrak{g}) / \mathfrak{m}_{\lambda}\) 是同构，并且 \(U(\mathfrak{g})\) 的每个有限余维双边理想恰好等于某个 \(m_{I}\)。
\item
  证明 \(\mathrm{U}(\mathfrak{g})\) 的主反自同构将 \(\mathfrak{m}_{\lambda}\) 变为 \(\mathfrak{m}_{\lambda^*}\)，其中 \(\lambda^{*} = -w_{0}\lambda\)（参见~命题11）。
\end{enumerate}

¶8) 令 g 为半单李代数；特别地，在本习题中不假定 g 是分裂的。令 \(\bar{k}\) 为 k 的代数闭包，并令

\[
\pi : \operatorname{Gal} (\bar {k} / k) \to \operatorname{Aut} (\bar {\mathrm{R}}, \bar {\mathrm{B}})
\]

为§5习题8中定义的同态。令 \(\operatorname{Gal}(\bar{k}/k)\) 通过 \(\pi\) 作用于 \(\bar{P}_{++}\)（\(\bar{R}\) 的相对于 \(\bar{B}\) 的支配权集合）；令 \(\Omega\) 为商 \(\bar{P}_{++}/\operatorname{Gal}(\bar{k}/k)\) 中元素的代表元系。

\begin{enumerate}
\def\labelenumi{\alph{enumi})}
\item
  置 \(\bar{g} = \bar{k} \otimes_{k} g\)。若 I 是 \(\bar{P}_{++}\) 的在 \(\operatorname{Gal}(\bar{k}/k)\) 下稳定的有限子集，则与 I 相关的双边理想 \(\bar{m}_{I}\) 是 \(\operatorname{U}(\bar{g})\) 的双边理想，形如 \(\bar{k} \otimes_{k} m_{I}\)，其中 \(m_{I}\) 是 \(\operatorname{U}(g)\) 的双边理想。证明 \(\operatorname{U}(g)\) 的每个有限余维双边理想都唯一地以这种方式得到。
\item
  令 \(\omega \in \Omega\)，令 \(\mathrm{I}(\omega)\) 为它在 \(\operatorname{Gal}(\bar{k}/k)\) 下的轨道，令 \(\mathrm{G}_{\omega}\) 为 \(\omega\) 的稳定子；令 \(k_{\omega}\) 为 \(\bar{k}\) 的子扩张，按 Galois 理论与 \(\mathrm{G}_{\omega}\) 对应。证明 \(\mathrm{U}(\mathfrak{g})/\mathfrak{m}_{\mathrm{I}(\omega)}\) 是中心同构于 \(k_{\omega}\) 的单代数。每个双边理想 \(\mathfrak{m}\)（\(\mathrm{U}(\mathfrak{g})\) 中的），若 \(\mathrm{U}(\mathfrak{g})/\mathfrak{m}\) 是有限维单代数，则  恰好是某个 \(\mathfrak{m}_{\mathrm{I}(\omega)}\)。
\item
  群 \(\operatorname{Gal}(\bar{k}/k)\) 通过 \(\pi\) 作用于表示环 \(\operatorname{R}(\mathfrak{g})\)；以 \(\operatorname{R}(\mathfrak{g})^{\mathrm{inv}}\) 表示 \(\operatorname{R}(\mathfrak{g})\) 中由 \(\operatorname{Gal}(\bar{k}/k)\) 不变元素组成的子环。证明映射 \([E] \to [\bar{k} \otimes_{k} E]\) 延拓为从 \(\operatorname{R}(\mathfrak{g})\) 到 \(\operatorname{R}(\mathfrak{g})^{\mathrm{inv}}\) 的单射同态，且其余核是挠群；该同态是同构，当且仅当对所有 \(\omega \in \Omega\)，\(\operatorname{U}(\mathfrak{g})/\mathfrak{m}_{\operatorname{I}(\omega)}\) 都是 \(k_{\omega}\) 上的矩阵代数；证明当 g 有 Borel 子代数时确实如此。\(^{14}\)
\end{enumerate}

\begin{enumerate}
\def\labelenumi{\arabic{enumi})}
\setcounter{enumi}{8}
\tightlist
\item
  令 \(\mathfrak{a}_1\) 和 \(\mathfrak{a}_2\) 为李代数。证明存在唯一的同态
\end{enumerate}

\[
f: \mathrm{R} (\mathfrak {a} _ {1}) \otimes_ {\mathbf {Z}} \mathrm{R} (\mathfrak {a} _ {2}) \to \mathrm{R} (\mathfrak {a} _ {1} \times \mathfrak {a} _ {2})
\]

使得 \(f([\mathrm{E}_1] \otimes [\mathrm{E}_2]) = [\mathrm{E}_1 \otimes \mathrm{E}_2]\)，当 \(\mathrm{E}_i\) 是 \(\mathfrak{a}_i\)-模（\(i = 1, 2\)）时。证明 \(f\) 是单射，并且当 \(\mathfrak{a}_1\) 和 \(\mathfrak{a}_2\) 为可分裂半单时为双射。

\begin{enumerate}
\def\labelenumi{\arabic{enumi})}
\setcounter{enumi}{9}
\tightlist
\item
  令 \(\Gamma\) 为 \(P\) 的包含 \(Q\) 的子群。这样的子群在 \(W\) 下稳定。
\end{enumerate}

\begin{enumerate}
\def\labelenumi{\alph{enumi})}
\item
  证明若 \(\lambda \in \mathrm{P}_{++} \cap \Gamma\)，则 \(\operatorname{E}(\lambda)\) 的每个权都属于 \(\Gamma\)。
\item
  令 \(\mathrm{R}_{\Gamma}(\mathfrak{g})\) 为 \(\mathrm{R}(\mathfrak{g})\) 的子群，其基为 \([\lambda]\)，其中 \(\lambda \in \mathrm{P}_{++} \cap \Gamma\)。若 \(\mathrm{E}\) 是 \(\mathfrak{g}\)-模，则 \([\mathrm{E}] \in \mathrm{R}_{\Gamma}(\mathfrak{g})\) 当且仅当 \(\mathrm{E}\) 的权属于 \(\Gamma\)。推出 \(\mathrm{R}_{\Gamma}(\mathfrak{g})\) 是 \(\mathrm{R}(\mathfrak{g})\) 的子环。
\item
  证明同态 \(\operatorname{ch} : \mathrm{R}_{\Gamma}(\mathfrak{g}) \to \mathbf{Z}[\Gamma]\) 是从 \(\mathrm{R}_{\Gamma}(\mathfrak{g})\) 到 \(\mathbf{Z}[\Gamma]\) 中由 W 不变元素组成的子环的同构。（使用定理2 (ii)。）
\end{enumerate}

\(d\) ) 具体描述 \(\mathrm{R}(\mathfrak{g})\) 和 \(\mathrm{R}_{\Gamma}(\mathfrak{g})\)：当 \(\mathfrak{g} = \mathfrak{sl}(2,k)\) 且 \(\Gamma = \mathrm{Q}\) 时。

¶11) 记号同第7节。对任意整数 \(m \geq 1\)，以 \(\Psi^{m}\) 表示 \(\mathbf{Z}[\Delta]\) 的自同态，它将 \(e^{\lambda}\) 变为 \(e^{m\lambda}\)。有 \(\Psi^{1} = \mathrm{Id}\) 且 \(\Psi^{m} \circ \Psi^{n} = \Psi^{mn}\)。

令 E 为有限维的 \(\Delta\)-分次向量空间。对所有 \(n \geq 0\)，以 \(a_{n}E\)（分别地，\(s_{n}E\)）表示 E 的第 n 个外幂（分别地，对称幂），并赋予其自然分次。

\begin{enumerate}
\def\labelenumi{\alph{enumi})}
\tightlist
\item
  证明
\end{enumerate}

\[
n \operatorname{ch} (s _ {n} \mathrm{E}) = \sum_ {m = 1} ^ {n} \Psi^ {m} (\operatorname{ch} (\mathrm{E})) \operatorname{ch} (s _ {n - m} \mathrm{E})
\]

以及

\[
n \operatorname{ch} (a _ {n} \mathrm{E}) = \sum_ {m = 1} ^ {n} (- 1) ^ {m - 1} \varPsi^ {m} (\operatorname{ch} (\mathrm{E})) \operatorname{ch} (a _ {n - m} \mathrm{E}).
\]

推出 \(\operatorname{ch}(s_{n}\operatorname{E})\) 和 \(\operatorname{ch}(a_{n}\operatorname{E})\) 可以表示为 \(\varPsi^{m}(\operatorname{ch}(\operatorname{E}))\)（\(1 \leq m \leq n\)）的有理系数多项式。例如：

\[
\mathrm{ch} (s _ {2} \mathrm{E}) = \frac {1}{2} \mathrm{ch} (\mathrm{E}) ^ {2} + \frac {1}{2} \varPsi^ {2} (\mathrm{ch} (\mathrm{E}))
\]

\[
\mathrm{ch} (a _ {2} \mathrm{E}) = \frac {1}{2} \mathrm{ch} (\mathrm{E}) ^ {2} - \frac {1}{2} \varPsi^ {2} (\mathrm{ch} (\mathrm{E})).
\]

\begin{enumerate}
\def\labelenumi{\alph{enumi})}
\setcounter{enumi}{1}
\tightlist
\item
  证明下列恒等式（在变量 T 的形式幂级数代数中，系数属于 \(Q[\Delta]\)）：
\end{enumerate}

\[
\sum_ {n = 0} ^ {\infty} \mathrm{ch} (s _ {n} \mathrm{E}) \mathrm{T} ^ {n} = \exp \left\{\sum_ {m = 1} ^ {\infty} \varPsi^ {m} (\mathrm{ch(E)}) \mathrm{T} ^ {m} / m \right\}
\]

以及

\[
\sum_ {n = 0} ^ {\infty} \mathrm{ch} (a _ {n} \mathrm{E}) \mathrm{T} ^ {n} = \exp \left\{\sum_ {m = 1} ^ {\infty} (- 1) ^ {m - 1} \varPsi^ {m} (\mathrm{ch(E)}) \mathrm{T} ^ {m} / m \right\}.
\]

\begin{enumerate}
\def\labelenumi{\alph{enumi})}
\setcounter{enumi}{2}
\tightlist
\item
  假设 \(\Delta\) 是一个群，从而 \(\varPsi^{m}\) 对所有 \(m\in \mathbf{Z}\) 都可以定义。证明若 \(\mathrm{E}^*\) 是 \(\mathrm{E}\) 的分次对偶，则
\end{enumerate}

\[
\operatorname{ch} \left(\mathrm{E} ^ {*}\right) = \Psi^ {- 1} (\operatorname{ch} (\mathrm{E})).
\]

\begin{enumerate}
\def\labelenumi{\alph{enumi})}
\setcounter{enumi}{3}
\tightlist
\item
  借助 ch 将 \(\mathrm{R}(\mathfrak{g})\) 与 Z[P] 的一个子环相识别。证明 \(\mathrm{R}(\mathfrak{g})\) 在 \(\varPsi^{m}\) 下稳定（\(m \in Z\)），并且前一习题中定义的子环 \(\mathrm{R}_{\varGamma}(\mathfrak{g})\) 也具有此性质。
\end{enumerate}

\begin{enumerate}
\def\labelenumi{\arabic{enumi})}
\setcounter{enumi}{11}
\tightlist
\item
  令 \(\lambda \in \mathrm{P}_{++}\)，令 \(\mathcal{X}_{\lambda}\) 为 \(\operatorname{E}(\lambda)\) 的权集。证明
\end{enumerate}

\[
\mathcal {X} _ {\lambda} \subset \lambda - \mathrm{P} _ {+ +}
\]

一般并不成立。（例如考虑 \(\mathfrak{sl}(3,k)\) 的伴随表示。）

\begin{enumerate}
\def\labelenumi{\arabic{enumi})}
\setcounter{enumi}{12}
\tightlist
\item
  令 \(\lambda = \sum_{\alpha \in B} a_{\alpha} \alpha\) 为 \(P_{++}\) 中的一个元素。对于 \(n = 0, 1, \ldots\)，以 \(\mathcal{X}_n\) 表示权 \(\mu\)，它是 \(E(\lambda)\) 的权，且满足 \(\lambda - \mu\) 是 n 个 \(n\)\(B\) 中元素之和。令 \(s_n\) 为 \(\mathcal{X}_n\) 中元素的重数之和（作为 \(E(\lambda)\) 的权的重数）。令 \(T = 2 \sum_{\alpha \in B} a_{\alpha}\)。证明：
\end{enumerate}

\begin{enumerate}
\def\labelenumi{\alph{enumi})}
\item
  T 是一个整数 \(\geq 0\)。
\item
  当 n > T 时，\(s_{n}=0\)，且 \(s_{T-n}=s_{n}\)。
\item
  若 r 是 T/2 的整数部分，则 \(s_{1} \leq s_{2} \leq \cdots \leq s_{r+1}\)。
\end{enumerate}

¶14) 令 \(\lambda \in \mathrm{P}_{++}\)，令 \(\mathrm{F}_{\lambda}\) 为 \(Z(\lambda)\) 的最大真子模，令 \(v\) 为 \(Z(\lambda)\) 中权为 \(\lambda\) 的本原元，参见~§6第3节。证明

\[
\mathrm{F} _ {\lambda} = \sum_ {\alpha \in \mathrm{B}} \mathrm{U} (\mathfrak {g}) X _ {- \alpha} ^ {\lambda (H _ {\alpha}) + 1} v = \sum_ {\alpha \in \mathrm{B}} \mathrm{U} (\mathfrak {n} _ {-}) X _ {- \alpha} ^ {\lambda (H _ {\alpha}) + 1} v.
\]

\begin{enumerate}
\def\labelenumi{\arabic{enumi})}
\setcounter{enumi}{14}
\tightlist
\item
  令 \(\lambda \in \mathfrak{h}^*\)，令 \(v\)（分别地，\(v'\)）为 \(\mathrm{Z}(\lambda)\)（分别地，\(\mathrm{E}(\lambda)\)）中权为 \(\lambda\) 的本原元。令 \(\mathrm{I}\)（分别地，\(\mathrm{I}'\)）为 \(v\)（分别地，\(v'\)）在 \(\mathrm{U}(\mathfrak{g})\) 中的湮没子。
\end{enumerate}

\[
a) \mathrm{I} = \mathrm{U} (\mathfrak {g}) \mathfrak {n} _ {+} + \sum_ {h \in \mathfrak {h}} \mathrm{U} (\mathfrak {g}) (h - \lambda (h)).
\]

\begin{enumerate}
\def\labelenumi{\alph{enumi})}
\setcounter{enumi}{1}
\item
I' 是不等于 U(g) 且包含 I 的 U(g) 的最大左理想。
\item
  若 \(\lambda \in \mathrm{P}_{++}\)，则
\end{enumerate}

\[
\mathrm{I} ^ {\prime} = \mathrm{I} + \sum_ {\alpha \in \mathrm{B}} \mathrm{U} (\mathfrak {g}) X _ {- \alpha} ^ {\lambda (H _ {\alpha}) + 1} = \mathrm{I} + \sum_ {\alpha \in \mathrm{B}} \mathrm{U} (\mathfrak {n} _ {-}) X _ {- \alpha} ^ {\lambda (H _ {\alpha}) + 1}.
\]

（使用前一习题。）

¶16) 令 V 和 V' 为两个 g-模。若存在一个线性映射 \(f : V \to V'\) 使其满足下列条件，则称 V' 从属于 V：

\(\alpha)\) \(f\) 是满射；\(\beta)\) \(f\) 下，\(\mathbf{V}\) 的本原元的像或者为 0，或者是 \(\mathrm{V}^{\prime}\) 的本原元；\(\gamma)\) \(f\) 是 \(\mathfrak{n}_{-}\)-同态。

\begin{enumerate}
\def\labelenumi{\alph{enumi})}
\item
  令 \(f: V \to V'\) 满足 \(\alpha\) )、\(\beta\) )、\(\gamma\) )。令 v 为 V 的本原元。则由 v 生成的 g-子模在 f 下的像是由 \(f(v)\) 生成的 g-子模。
\item
  令 \(f: \mathrm{V} \to \mathrm{V}'\) 满足 \(\alpha), \beta), \gamma)\)。令 \(\mathrm{W}\) 为 \(\mathfrak{g}\) 的 \(\mathrm{V}\)-子模。则 \(f(\mathrm{W})\) 是 \(\mathfrak{g}\) 的 \(\mathrm{V}'\)-子模，并且从属于 \(\mathrm{W}\)。
  \item
  令 \(V = E_{1} \oplus \cdots \oplus E_{s}\) 且 \(V' = E_{1}' \oplus \cdots E_{s'}\) 为 V 和 \(V'\) 分解为单模之和的分解。则 \(V'\) 从属于 V，当且仅当 \(s' \leq s\)，且存在 \(\sigma \in S_{s}\)，使得 \(E_{i}'\) 从属于 \(\mathrm{E}_{\sigma(i)}\)，对 \(i = 1, \ldots, s'\) 成立。
\item
  若 \(V'\) 从属于 V 且 V 是单的，则 \(V'\) 是单的或化为 0。
\item
  假设 V 和 \(V'\) 是单的。令 \(\lambda\) 和 \(\lambda'\) 为它们的最高权。则 \(V'\) 从属于 V，当且仅当 \(\lambda'(H_{\alpha}) \leq \lambda(H_{\alpha})\) 对所有 \(\alpha \in B\) 成立。（充分性使用前一习题。）
\end{enumerate}

\begin{enumerate}
\def\labelenumi{\arabic{enumi})}
\setcounter{enumi}{16}
\tightlist
\item
  令 \(\lambda, \mu \in \mathrm{P}_{++}\) 和 \(\alpha \in \mathrm{B}\) 满足 \(\lambda(H_{\alpha}) \geq 1\) 且 \(\mu(H_{\alpha}) \geq 1\)。令 \(\mathrm{F} = \mathrm{E}(\lambda) \otimes \mathrm{E}(\mu)\)。
\end{enumerate}

\begin{enumerate}
\def\labelenumi{\alph{enumi})}
\item
  证明 \(\dim \mathrm{F}^{\lambda +\mu} = 1\)，且 \(\dim \mathrm{F}^{\lambda +\mu -\alpha} = 2\)。
\item
  证明 \(X_{\alpha}: F^{\lambda+\mu-\alpha} \to F^{\lambda+\mu}\) 是满射，且其核中的非零元素是本原元（注意，若 \(\beta \in B\) 异于 \(\alpha\)，则 \(\lambda + \mu - \alpha + \beta\) 不是 F 的权）。
\item
  推出 \(\mathrm{E}(\lambda)\otimes\mathrm{E}(\mu)\) 含有唯一一个同构于
\end{enumerate}

\[
\mathrm{E} (\lambda + \mu - \alpha).
\]

\begin{enumerate}
\def\labelenumi{\alph{enumi})}
\setcounter{enumi}{3}
\tightlist
\item
  证明 \(\mathbf{S}^{2}(\mathrm{E}(\lambda))\)（分别地，\(\bigwedge^{2}\mathrm{E}(\lambda)\)）含有唯一一个同构于 \(\mathrm{E}(2\lambda)\)（分别地，同构于 \(\mathrm{E}(2\lambda-\alpha)\)）的子模。
\end{enumerate}

¶ 18) 选取一个正定非退化对称双线性形式 \((\cdot|\cdot)\) 在 \(\mathfrak{h}_{\mathbf{R}}^{*}\) 上 W-不变。令 \(\lambda \in P_{++}\)。

\begin{enumerate}
\def\labelenumi{\alph{enumi})}
\item
  令 \(\mu\) 为 \(\mathrm{E}(\lambda)\) 的一个权。将 \(\lambda - \mu\) 写成 \(\sum_{\alpha \in \mathrm{B}} k_{\alpha} \alpha\) 的形式。令 \(\alpha \in \mathrm{B}\) 满足 \(k_{\alpha} \neq 0\)。证明存在 \(\alpha_{1}, \ldots, \alpha_{n} \in \mathrm{B}\)，使得 \((\lambda | \alpha_{1}) \neq 0\)、\((\alpha_{1} | \alpha_{2}) \neq 0, \ldots, (\alpha_{n-1} | \alpha_{n}) \neq 0\)、\((\alpha_{n} | \alpha) \neq 0\)。
\item
  令 \(v\) 为 \(\mathrm{E}(\lambda)\) 的本原元，并令 \(\alpha_1, \ldots, \alpha_n \in \mathrm{B}\) 满足下列条件：
\end{enumerate}

\begin{enumerate}
\def\labelenumi{(\roman{enumi})}
\item
  \((\alpha_{i}|\alpha_{i+1})\neq0\) 对 \(i=1,2,\ldots,n-1;\)
\item
  \((\alpha_{i}|\alpha_{j}) = 0\) 对 \(j > i + 1\)；
\item
  \(\lambda(H_{\alpha_{1}})\neq0\) 且 \(\lambda(H_{\alpha_{2}})=\cdots=\lambda(H_{\alpha_{n}})=0.\)
\end{enumerate}

证明 \(X_{-\alpha_n}X_{-\alpha_{n-1}} \ldots X_{-\alpha_1}v \neq 0\)。(注意，对于 \(1 \leq s \leq n\)，

\[
\lambda - \alpha_ {1} - \dots - \alpha_ {s - 1} + \alpha_ {n}
\]

不是 \(\mathrm{E}(\lambda)\) 的权，并由归纳法（以 \(X_{\alpha_s}X_{-\alpha_s}X_{-\alpha_{s-1}}\ldots X_{-\alpha_1}v\neq 0\) 为归纳变量）推出 \(s\)。) 若 \(\sigma \in \mathfrak{S}_n\) 且 \(\sigma \neq 1\)，则 \(X_{-\alpha_{\sigma(n)}}X_{-\alpha_{\sigma(n-1)}}\ldots X_{-\alpha_{\sigma(1)}}v = 0\)。(令 \(r\) 为满足 \(\sigma(r)\neq r\) 的最小整数。利用 \(a\) ) 证明 \(\lambda -\alpha_{\sigma(1)} - \ldots -\alpha_{\sigma(r)}\) 不是 \(\mathrm{E}(\lambda)\) 的权。)

\begin{enumerate}
\def\labelenumi{\alph{enumi})}
\setcounter{enumi}{2}
\tightlist
\item
  令 \(\lambda' \in \mathrm{P}_{++}\)。连接 \(\lambda\) 与 \(\lambda'\) 的链是 B 中元素的序列 \((\alpha_1, \ldots, \alpha_n)\)，满足 \(n \geq 1\)、\((\lambda|\alpha_1) \neq 0\)、\((\alpha_1|\alpha_2) \neq 0, \ldots, (\alpha_{n-1}|\alpha_n) \neq 0\)、\((\alpha_n|\lambda') \neq 0\)。若这样的链的 \((\alpha_1, \ldots, \alpha_n)\) 没有连接 \(\lambda\) 与 \(\lambda'\) 的真子序列，则称其为极小链。在此情形，有：\((\alpha_i|\alpha_j) = 0\)（若 \(|i - j| \geq 2\)）；若 \((\lambda|\alpha_i) = 0\)，则 \(i \geq 2\)；若 \((\lambda'|\alpha_i) = 0\)，则 \(i \leq n - 1\)。
\end{enumerate}

\(d\) ) 令 \((\alpha_{1},\ldots ,\alpha_{n})\) 为连接 \(\lambda\) 与 \(\lambda^{\prime}\) 的极小链。若 \(v^{\prime}\) 是 \(\operatorname {E}(\lambda ')\) 中的本原向量，置：

\[
v _ {s} = X _ {- \alpha_ {s}} X _ {- \alpha_ {s - 1}} \ldots X _ {- \alpha_ {1}} v (s = 0, 1, \ldots , n)
\]

\[
v _ {s} ^ {\prime} = X _ {- \alpha_ {s + 1}} X _ {- \alpha_ {s + 2}} \dots X _ {- \alpha_ {n}} v ^ {\prime} (s = 0, 1, \ldots , n)
\]

\(a_0 = (\lambda |\alpha_1), a_n = (-1)^n (\lambda '|\alpha_n), a_s = (-1)^{s + 1}(\alpha_s|\alpha_{s + 1}), 1 \leq s \leq n - 1.\) 证明

\[
\sum_ {s = 0} ^ {n} a _ {s} v _ {s} \otimes v _ {s} ^ {\prime}
\]

是 \(\mathrm{E}(\lambda)\otimes\mathrm{E}(\lambda')\) 中权为 \(\lambda+\lambda'-\alpha_{1}-\cdots-\alpha_{n}\) 的本原元，并且在相差一个伸缩的意义下唯一。

（使用 b) 和 c) 证明 \(\mathrm{E}(\lambda) \otimes \mathrm{E}(\lambda')\) 中每个权为

\[
\lambda + \lambda^ {\prime} - \alpha_ {1} - \dots - \alpha_ {n}
\]

的是 \(v_{s} \otimes v_{s}^{\prime}\) 的线性组合。然后写出这种线性组合成为本原向量的条件。）

推出 \(\mathrm{E}(\lambda) \otimes \mathrm{E}(\lambda')\) 含有唯一一个同构于

\[
\operatorname{E} \left(\lambda + \lambda^ {\prime} - \alpha_ {1} - \dots - \alpha_ {n}\right).
\]

（当 \(n = 1\) 时，我们得到习题17。）

\begin{enumerate}
\def\labelenumi{\alph{enumi})}
\setcounter{enumi}{4}
\tightlist
\item
  令 w 为 \(\mathrm{E}(\lambda)\otimes\mathrm{E}(\lambda')\) 的本原元。假设 w 的权 \(\nu\) 异于 \(\lambda+\lambda'\)。证明存在一个 \((\alpha_{1},\ldots,\alpha_{n})\)，它连接 \(\lambda\) 与 \(\lambda'\)，使得
\end{enumerate}

\[
\nu \leq \lambda + \lambda^ {\prime} - \alpha_ {1} - \dots - \alpha_ {n}.
\]

(令 C 为由 \(\alpha \in \mathrm{B}\) 中满足 \(\alpha\) 在 \(\lambda + \lambda' - \nu\) 中的指标为 \(\neq 0\) 的元素组成的集合。令 D（分别地，D'）为满足下述条件的 \(\alpha \in \mathrm{C}\) 的集合：存在 \(\gamma_1, \ldots, \gamma_t \in \mathrm{C}\)，使得 \((\lambda|\gamma_1) \neq 0\)（分别地，\((\lambda'|\gamma_1) \neq 0\)）、\((\gamma_1|\gamma_2) \neq 0, \ldots, (\gamma_{t-1}|\gamma_t) \neq 0\)、\((\gamma_t|\alpha) \neq 0\)。令 Y（分别地，Y'）为 E( \(\lambda\) )（分别地，E( \(\lambda'\) )）中形如 \(\lambda - \sum_{\alpha \in \mathrm{C}} k_\alpha\alpha\)（分别地，\(\lambda' - \sum_{\alpha \in \mathrm{C}} k_\alpha\alpha\)）且 \(k_\alpha \in \mathbf{N}\) 的权的集合。证明 \(w\) 属于 \(\left(\sum_{\mu \in \mathrm{Y}}\mathrm{E}(\lambda)^{\mu}\right)\otimes \left(\sum_{\mu^{\prime}\in \mathrm{Y}^{\prime}}\mathrm{E}(\lambda^{\prime})^{\mu^{\prime}}\right)\)。利用 \(a\) ) 以及 \(\nu \neq \lambda +\lambda '\) 这一事实，证明 \(\mathrm{D}\cap \mathrm{D}'\neq \emptyset .\))

\(f\) ) 证明下列性质等价：

\begin{enumerate}
\def\labelenumi{(\roman{enumi})}
\item
  \(\mathrm{E}(\lambda)\otimes \mathrm{E}(\lambda^{\prime})\) 同构于 \(\mathrm{E}(\lambda +\lambda^{\prime})\)
\item
  \(\mathrm{E}(\lambda) \otimes \mathrm{E}(\lambda')\) 是单模。
\item
  不存在连接 \(\lambda\) 与 \(\lambda'\) 的链。
\item
  R 是两个根系 \(R_{1}\) 和 \(R_{1}^{\prime}\) 的直和，使得 \(\lambda \in \mathrm{P}(\mathrm{R}_{1})\) 且 \(\lambda^{\prime} \in \mathrm{P}(\mathrm{R}_{1}^{\prime})\)。
\item
  \(\mathfrak{g}\) 是两个理想 \(\mathfrak{s}\) 和 \(\mathfrak{s}'\) 的直积，使得 \(\mathfrak{s}'\) .E(λ) = 0 且 \(\mathfrak{s}\) .E(λ') = 0。
\end{enumerate}

（使用 \(d\) ) 证明 (ii) 与 (iii) 等价。）\(^{15}\)

\begin{enumerate}
\def\labelenumi{\arabic{enumi})}
\setcounter{enumi}{18}
\item
  回忆命题10的记号。令 \(\bar{k}\) 为 \(k\) 的代数闭包。若 \(x \in \mathfrak{g}\)，以 \(\mathcal{X}_{\mathrm{E}}(x)\)（分别地，\(\mathcal{X}_{\mathrm{F}}(x)\)、分别地，\(\mathcal{X}_{\mathrm{G}}(x)\)）表示 \(x_{\mathrm{E}}\)（分别地，\(x_{\mathrm{F}}\)、分别地，\(x_{\mathrm{G}}\)）在 \(\bar{k}\) 中的特征值集合。证明 \(\mathcal{X}_{\mathrm{G}}(x) = \mathcal{X}_{\mathrm{E}}(x) + \mathcal{X}_{\mathrm{F}}(x)\) 对所有 \(x \in \mathfrak{g}\) 成立，并且在给定 E 和 F 时，此性质在同构意义下刻画单 \(\mathfrak{g}\)-模 G。
\item
  假设 \(k = \mathbf{R}\) 或 \(\mathbf{C}\)。令 \(\Gamma\) 为李代数为 \(\mathfrak{g}\) 的单连通李群。令 \(\lambda, \mu, \mathrm{E}, \mathrm{F}, \mathrm{G}\) 如命题10中。令 \((e_1, \ldots, e_n)\)（分别地，\((f_1, \ldots, f_p)\)）为 \(\mathrm{E}\)（分别地，\(\mathrm{F}\)）的一个由 \(\mathfrak{h}\) 的特征向量组成的基，其中 \(e_1 \in \mathrm{E}^\lambda\) 且 \(f_1 \in \mathrm{F}^\mu\)。可将 \(\mathrm{E}, \mathrm{F}, \mathrm{G}\) 看作 \(\Gamma\)-模。若 \(\gamma \in \Gamma\)，以 \(a_i(\gamma)\) 表示 \(\gamma.e_1\) 的指标为 \(i\) 的坐标，以 \(b_j(\gamma)\) 表示 \(\gamma.f_1\) 的指标为 \(j\) 的坐标。证明 \(a_i b_j\) 上的函数 \(\Gamma\) 不恒为零。推出对所有 \((i,j)\)，存在 \(\mathrm{G} \subset \mathrm{E} \otimes \mathrm{F}\) 中的元素，其指标为 \((i,j)\) 的坐标 \(\neq 0\)；当 \(k = \mathbf{R}\) 或 \(\mathbf{C}\) 时，这给出命题10的一个新证明。由此转到 \(k = \mathbf{Q}\) 的情形，再转到任意域的情形，参见~习题4。
\item
  令 \(\lambda, \mu \in \mathrm{P}_{++}\)。令 E、F、G 为单 \(\mathfrak{g}\)-模，其最高权为 \(\lambda, \mu, \lambda + \mu\)，令 \(n\) 为满足 \(\geq 1\) 的整数。若 \(\omega\) 是 E 的一个重数为 \(n\) 的权，则 \(\omega + \mu\) 是 G 的一个重数 \(\geq n\) 的权。（有 \(\mathrm{G}^{\omega + \mu} \subset \bigoplus_{\nu + \sigma = \omega + \mu} \mathrm{E}^{\nu} \otimes \mathrm{F}^{\sigma}\)。若 \(\dim \mathrm{G}^{\omega + \mu} < n\)，则 \(\mathrm{G}^{\omega + \mu}\) 到 \(\mathrm{E}^{\omega} \otimes \mathrm{F}^{\mu}\) 的投影形如 \(\mathrm{E}' \otimes \mathrm{F}^{\mu}\)，其中 \(\mathrm{E}'\) 真包含于 \(\mathrm{E}^{\omega}\)。通过选取 E 和 F 的适配基并仿照命题10的证明，由此导出矛盾。）
\item
  假设 \(\mathfrak{g}\) 是单的，换言之，R 是不可约的。证明存在唯一的支配权 \(\lambda \neq 0\)，使得 \(\mathrm{E}(\lambda)\) 的权集为 \(\mathrm{W}.\lambda \cup \{0\}\)：有 \(\lambda = \alpha\)，其中 \(\alpha \in \mathrm{R}\) 满足 \(H_{\alpha}\) 是 \(\mathrm{R}^{\vee}\) 的最高根。当所有根等长时（情形 \(\mathrm{A}_l,\mathrm{D}_l,\mathrm{E}_6,\mathrm{E}_7,\mathrm{E}_8\)），有 \(\lambda = \tilde{\alpha}\)；这是唯一一个同时为支配权的根；相应的表示是 g 的伴随表示。在其他情形中，\(\lambda\) 是长度最小且为支配权的唯一根；按第VI章图版中的记号，有：\(\lambda = \varpi_{1}\)（\(B_{l}\) 型）；\(\lambda = \varpi_{2}\)（\(C_{l}\) 型）；\(\lambda = \varpi_{4}\)（\(F_{4}\) 型）；\(\lambda = \varpi_{1}\)（\(G_{2}\) 型）。
\item
  令 \(\mathrm{U}^0\) 为 \(\mathfrak{h}\) 在 \(\mathrm{U}(\mathfrak{g})\) 中的交换子（参见~§6第4节）。
\end{enumerate}

\begin{enumerate}
\def\labelenumi{\alph{enumi})}
\item
  令 V 为一个（有限维）g-模。证明各 \(V^{\lambda}\)（\(\lambda \in P\)）在 \(U^{0}\) 下稳定，并且若 V 是单的且 \(V^{\lambda} \neq 0\)，则 \(V^{\lambda}\) 是单 \(U^{0}\)-模（使用 \(\mathrm{U}(\mathfrak{g}) = \oplus\mathrm{U}^{\lambda}\) 的分解，同上）。
\item
  证明对于每个 \(c \neq 0\)（属于 \(U^{0}\)），都存在 \(\rho\) 的有限维单表示 \(U^{0}\)，使得 \(\rho(c) \neq 0\)。（使用 a) 和第I章§7的习题3 a)。
\end{enumerate}

\begin{enumerate}
\def\labelenumi{\arabic{enumi})}
\setcounter{enumi}{23}
\tightlist
\item
  令 \(A\) 为幺元结合代数，\(M\) 为其有限余维双边理想的集合。
\end{enumerate}

\begin{enumerate}
\def\labelenumi{\alph{enumi})}
\tightlist
\item
  令 \(\mathrm{U}^*\) 为 U 的向量空间对偶。若 \(\theta \in \mathrm{U}^*\)，证明下列性质等价：
\end{enumerate}

\begin{enumerate}
\def\labelenumi{(\roman{enumi})}
\item
  存在 \(\mathfrak{m} \in \mathrm{M}\)，使得 \(\theta(\mathfrak{m}) = 0\)；
\item
   存在两个有限族 \((\theta_{i}^{\prime})\) 和 \((\theta_{i}^{\prime\prime})\)，它们由 \(U^{*}\) 中的元素组成，使得
\end{enumerate}

\[
\theta (x y) = \sum_ {i} \theta_ {i} ^ {\prime} (x) \theta_ {i} ^ {\prime \prime} (y) \quad \text { for   all } x, y \in \mathrm{U}.
\]

具有这些性质的元素 \(\theta\) 构成 \(\mathrm{U}'\) 的子空间，它与《代数》第III章§11习题27中记为 \(\mathrm{U}^*\)\(\mathrm{B}'\) 的空间重合。\(\mathrm{U}'\) 上存在唯一的余代数结构，其余积 \(c: \mathrm{U}' \to \mathrm{U}' \otimes \mathrm{U}'\) 由下式给出

\[
c (\theta) = \sum_ {i} \theta_ {i} ^ {\prime} \otimes \theta_ {i} ^ {\prime \prime},
\]

其中 \(\theta_i', \theta_i'' \in \mathrm{U}'\) 满足 \(\theta(xy) = \sum_{i} \theta_i'(x) \theta_i''(y)\) 对所有 \(x, y \in \mathrm{U}\) 成立，参见~(ii)。

余代数 \(\mathrm{U}'\) 是由子空间 \((\mathrm{U} / \mathfrak{m})^*\)（\(\mathfrak{m} \in \mathrm{M}\)）构成的递增滤过的并；这些子空间都是有限维的。\(\mathrm{U}'\) 的对偶可以与代数 \(\hat{\mathrm{U}} = \varprojlim \mathrm{U} / \mathfrak{m}\) 相识别；若给 \(\hat{\mathrm{U}}\) 赋予 \(\mathrm{U} / \mathfrak{m}, \mathfrak{m} \in \mathrm{M}\) 上的离散拓扑的射影极限拓扑，则 \(\hat{\mathrm{U}}\) 上的连续线性形式正是 \(\mathrm{U}'\) 中的元素。


\begin{enumerate}
\def\labelenumi{\alph{enumi})}
\setcounter{enumi}{1}
\tightlist
\item
  令 E 为有限维左 U-模。其湮没子 \(m_{E}\) 属于 M；复合 \(\hat{U} \to U/m_{E} \to \text{End}(E)\) 赋予 E 一个左 \(\hat{U}\)-模结构。若 F 为有限维左 U-模，则线性映射 \(f : E \to F\) 是 U-同态，当且仅当它是 \(\hat{U}\)-同态。若 \(a \in E, b \in E^{*}\)，线性形式 \(\theta_{a,b} : x \mapsto \langle xa, b \rangle\) 属于 \(U'\)，且
\end{enumerate}

\[
\langle x, \theta_ {a, b} \rangle = \langle x a, b \rangle \quad \text { for   all } x \in \hat {\mathrm{U}}.
\]

  当 E 变化时，\(\theta_{a,b}\)（其中 \(a,b\) 变化）生成 \(k\)-向量空间 U'。

\begin{enumerate}
\def\labelenumi{\alph{enumi})}
\setcounter{enumi}{2}
\tightlist
\item
  令 \(X_{U}\) 为有限维左 U-模的同构类集合。对所有 \(E \in X_{U}\)，令 \(u_{E}\) 为 E 的 k-线性自同态；假设
\end{enumerate}

\[
f \circ u _ {\mathrm{E}} = u _ {\mathrm{F}} \circ f \quad \text { for   all   E, F } \in X_{U} \text { and } f \in \operatorname{Hom} _ {\mathrm{U}} (\mathrm{E,F}).
\]

证明存在唯一的元素 \(x \in \hat{U}\)，使得 \(x_{E} = u_{E}\) 对所有 \(E \in X_{U}\) 成立。（化归到 U 为有限维的情形。）

¶ 25) 令 a 为李代数，U 为其包络代数。将习题24的定义和结果应用于代数 U。特别地，\(U' \subset U^*\)，且 \(U'\) 的对偶可以与代数 \(\hat{U} = \varprojlim U/m\) 相识别；典则映射 \(U \to \hat{U}\) 是单射（第I章§7的习题3）。

\begin{enumerate}
\def\labelenumi{\alph{enumi})}
\tightlist
\item
  U 的余代数结构（第II章§1第4节）在对偶 U* 上定义一个代数结构（参见~第II章§1第5节的命题10以及《代数》第III章§11第2节）。若 E、F 是有限维 a-模，
\end{enumerate}

\[
\theta_ {a, b} \circ \theta_ {c, d} = \theta_ {a \otimes c, b \otimes d} \text { for } a \in \mathrm{E}, b \in \mathrm{E} ^ {*}, c \in \mathrm{F}, d \in \mathrm{F} ^ {*}.
\]

推出 \(U'\) 是 \(U^{*}\) 的子代数。\(U'\) 的余代数结构和代数结构使其成为交换双代数（《代数》第III章§11第4节）。

\begin{enumerate}
\def\labelenumi{\alph{enumi})}
\setcounter{enumi}{1}
\tightlist
\item
  令 \(x\) 为 \(\hat{\mathrm{U}}\) 的一个元素；将 \(x\) 与线性形式 \(\mathrm{U}' \to k\) 相识别。证明下列性质等价：
\end{enumerate}

\begin{enumerate}
\def\labelenumi{(\roman{enumi})}
\item
  \(x\) 是从 \(\mathrm{U}'\) 到 \(k\) 的代数同态。
\item
  对所有有限维 a-模 E 和 F，都有 \(x_{E \otimes F} = x_{E} \otimes x_{F}\)。
\end{enumerate}

（先证明 (ii) 等价于

\[
\left(\mathrm{ii} ^ {\prime}\right) x \left(\theta_ {a \otimes c, b \otimes d}\right) = x \left(\theta_ {a, b}\right) x \left(\theta_ {c, d}\right) \text {if} a \in \mathrm{E}, b \in \mathrm{E} ^ {*}, c \in \mathrm{F}, d \in \mathrm{F} ^ {*},
\]

并使用 \(\theta_{a,b}\) 生成 \(\mathrm{U}'\) 这一事实。）

\(c)\) 令 \(x\) 为满足 \(\hat{\mathbf{U}}\) ) 中条件 (i) 和 (ii) 的 \(b\) 元素。证明下列条件等价：

\begin{enumerate}
\def\labelenumi{(\roman{enumi})}
\setcounter{enumi}{2}
\item
  \(x\) 将 \(\mathrm{U}'\) 的单位元映为 \(k\) 的单位元。
\item
  \(x \neq 0.\)
\item
  若给 \(k\) 赋予平凡的 \(\mathfrak{a}\)-模结构，则有 \(x_{k} = \mathrm{Id}\)。
\item
  对所有 E，\(x_{\mathrm{E}}\) 都是可逆的。
\end{enumerate}

（(iii) \(\Longleftrightarrow\) (iv) 的等价性由 \(x\) 是代数同态这一事实得到。另一方面，\(x_{k} = \lambda \mathrm{Id}\)，其中 \(\lambda \in k\)。利用 \(\mathfrak{a}\)-同构 \(k \otimes \mathrm{E} \to \mathrm{E}\)，推出对所有 E 都有 \(\lambda x_{\mathrm{E}} = x_{\mathrm{E}}\)，特别地，取 \(\lambda^2 = \lambda\) 得 \(\mathrm{E} = k\)。\(\lambda = 1\) 的情形对应于 \(x \neq 0\)，因而 (iv) \(\Longleftrightarrow\) (v)，且 (vi) \(\Longrightarrow\) (v)。为证明 (v) \(\Longrightarrow\) (vi)，证明若 F 是 E 的对偶，则 \(^t x_{\mathrm{F}} \circ x_{\mathrm{E}} = \lambda \mathrm{Id}_{\mathrm{E}}\)。）

\begin{enumerate}
\def\labelenumi{\alph{enumi})}
\setcounter{enumi}{3}
\tightlist
\item
  令 G 为 \(\hat{U}\) 中满足上述条件 (i) 至 (vi) 的元素集。证明 G 是 \(\hat{U}\) 的可逆元素群的子群。
\end{enumerate}

令 \(x \in G\)。若 E 是有限维 \(\mathfrak{a}\)-模，则 \(x_{\mathrm{E}} \in \mathbf{GL}(\mathrm{E})\)。特别地，取 \(\mathrm{E} = \mathfrak{a}\)，赋予其伴随表示；这给出元素 \(x_{\mathfrak{a}} \in \mathbf{GL}(\mathfrak{a})\)。证明 \(x_{\mathfrak{a}}\) 是 \(\mathfrak{a}\) 的自同构（使用由括号给出的 \(\mathfrak{a}\)-同态 \(\mathfrak{a} \otimes \mathfrak{a} \to \mathfrak{a}\)）。这给出同态 \(v: \mathrm{G} \to \operatorname{Aut}(\mathfrak{a})\)。若 \(\mathrm{E}\) 是有限维 \(\mathfrak{a}\)-模，

\[
x _ {\mathrm{E}} (y. e) = v (x) (y). x _ {\mathrm{E}} (e) \quad \mathrm{if} y \in \mathfrak {a}, e \in \mathrm{E}
\]

（使用由 \(\mathfrak{a}\) 在 \(\mathfrak{a} \otimes \mathrm{E} \to \mathrm{E}\) 上的作用给出的 \(\mathfrak{a}\)-同态 \(\mathrm{E}\)。）

\begin{enumerate}
\def\labelenumi{\alph{enumi})}
\setcounter{enumi}{4}
\tightlist
\item
  U 的主反自同构 \(\sigma\) 由连续性延拓到 \(\hat{\mathbf{U}}\)。其转置保持 \(\mathrm{U}'\) 稳定，并在 \(\mathrm{U}'\) 上诱导逆（《代数》第III章§11习题4）。若 \(x \in \mathrm{G}\)，则 \(\sigma(x) = x^{-1}\)。
\end{enumerate}

\(f\) ) 假设 \(\mathfrak{a}\) 是半单的 \(^{16}\)。令 \(n\) 为 \(\mathfrak{a}\) 的幂零元。则存在唯一的 \(e^n\) 中的元素 \(\mathrm{G}\)，使得对每个有限维 \((e^n)_{\mathrm{E}} = \exp (n_{\mathrm{E}})\)-模 E 都有 \(\mathfrak{a}\)。有 \(v(e^n) = \exp (\operatorname {ad}n)\in \operatorname {Aut}(\mathfrak{a})\)，所以 \(\operatorname {Aut}_e(\mathfrak{a})\subset v(\mathrm{G})\)。

证明若 b 是由幂零元素组成的 a 的子代数，则

\[
e ^ {n}. e ^ {m} = e ^ {\mathrm{H} (n, m)} \quad \text { for } n, m \in \mathfrak {b},
\]

其中 H 表示 Hausdorff 级数（第II章§6）。

¶ 26) 将习题25的记号和结果应用于 a = g 的情形（分裂情形）。

\begin{enumerate}
\def\labelenumi{\alph{enumi})}
\item
  令 \(x \in G\)，令 \(\sigma = v(x)\) 为它在 \(\operatorname{Aut}(\mathfrak{g})\) 中的像。若 \(\rho\) 是 g 的一个表示，则 \(\rho\) 与 \(\rho \circ \sigma\) 等价。推出（参见~第2节，注1）\(\sigma\) 属于 \(\operatorname{Aut}_{0}(\mathfrak{g})\)。将此结果推广到任意半单代数。
\item
  令 \(\varphi \in \mathrm{T}_{\mathrm{P}} = \mathrm{Hom}(\mathrm{P}, k^{*})\)，其中 \(\mathrm{P} = \mathrm{P}(\mathrm{R})\)。若 \(\mathrm{E}\) 是 \(\mathfrak{g}\)-模，令 \(\varphi_{\mathrm{E}}\) 为 \(\mathrm{E}\) 的自同态，其在每个 \(\mathrm{E}^{\lambda}\)（\(\lambda \in \mathrm{P}\)）上的限制是比值为 \(\varphi(\lambda)\) 的伸缩。证明存在唯一的元素 \(t(\varphi) \in \mathrm{G}\)，使得对所有 \(t(\varphi)_{\mathrm{E}} = \varphi_{\mathrm{E}}\) 都有 \(\mathrm{E}\)（使用习题24 c) 以及习题25的刻画 (ii) 和 (vi)）。这给出同态 \(t: \mathrm{T}_{\mathrm{P}} \to \mathrm{G}\)。证明 \(t\) 是单射。利用这一点将 \(\mathrm{T}_{\mathrm{P}}\) 与 \(\mathrm{G}\) 的一个子群相识别。复合映射 \(\mathrm{T}_{\mathrm{P}} \to \mathrm{G} \to \mathrm{Aut}(\mathfrak{g})\) 是§5第2节中记为 \(f \circ q\) 的同态。
\item
  令 \(x \in G\) 满足 \(\sigma = v(x)\) 属于 \(f(\mathrm{T}_{\mathrm{Q}})\) 的子群 \(\mathrm{Aut}_{0}(\mathfrak{g})\)（§5第2节），换言之，它在 h 上作用平凡；以 \(T_{Q} = \mathrm{Hom}(Q, k^{*})\) 表示与 x 对应的 \(\psi\) 中的元素。我们将证明 x 属于 \(T_{P}\)。依次证明：
\end{enumerate}

\(c_{1}\) ) 若 E 是 \(\mathfrak{g}\)-模，则 \(x_{\mathrm{E}}\) 是 E 的 \(\mathfrak{h}\)-自同态。

（使用 \(g \otimes E \to E\) 的 g-同态，以及 x 在 h 上作用平凡这一事实。）特别地，各 \(E^{\mu}\) 在 \(x_{E}\) 下稳定。

\(c_{2}\) ) 存在 \(\varphi \in \mathrm{T}_{\mathrm{P}}\)，使得对于每个 \(\mathfrak{g}\)-模 E 以及 E 中每个权为 \(e\) 的本原元 \(\lambda\)，都有 \(x_{\mathrm{E}}e = \varphi (\lambda)e\)。

（选取 \(\varphi\)，使得当 E 是基本模 \(\mathrm{E}(\varpi_{\alpha})\) 时该关系成立。利用将这类模嵌入若干 \(\mathrm{E}(\lambda)\) 的张量积，推出 \(\lambda \in \mathrm{P}_{++}\)、\(\mathrm{E}(\varpi_{\alpha})\) 的情形。再由此推广到一般情形。）

\(c_{3}\) ) 按 \(\varphi\) ) 选取 \(c_{2}\)。令 E 为最高权为 \(\lambda\) 的单 g-模，令 \(\mu\) 为 E 的一个权；则 \(\lambda - \mu \in Q\)。证明 \(x_{E}\) 在 \(E^{\mu}\) 上的限制是比值为 \(\varphi(\lambda)\psi(\mu - \lambda)\) 的伸缩。（方法同 \(c_{1}\) )。）

\(c_{4}\) ) 若 \(\lambda, \mu \in P_{++}\)，且 \(\alpha \in B\) 不与 \(\lambda\) 或 \(\mu\) 中任一个正交，则

\[
\varphi (\lambda + \mu - \alpha) = \varphi (\lambda + \mu) \psi (- \alpha),
\]

所以 \(\varphi (\alpha) = \psi (\alpha)\)。（使用 \(c_{2}\) )、\(c_{3}\) ) 以及 \(\operatorname {E}(\lambda +\mu -\alpha)\) 到 \(\operatorname {E}(\lambda)\otimes \operatorname {E}(\mu)\) 中的嵌入，参见~习题17。）

\(c_{5})\) 由 \(c_{4})\) 推出 \(\varphi|\mathbf{Q}=\psi\)，并使用 \(c_{3})\) 推出 \(x=t(\varphi)\)。\(d)\) 通过 \(\mathrm{T}_{\mathrm{Q}}\) 将 \(\mathrm{Aut}_{0}(\mathfrak{g})\) 与 \(f\) 的一个子群相识别。由 \(a)\)，有

\[
\operatorname{Aut} _ {e} (\mathfrak {g}) \subset v (\mathrm{G}) \subset \operatorname{Aut} _ {0} (\mathfrak {g})
\]

且由 \(c\) )，\(v(\mathrm{G}) \cap \mathrm{T}_{\mathrm{Q}} = \mathrm{Im}(\mathrm{T}_{\mathrm{P}})\)。推出（参见~§5第3节）\(\mathrm{Aut}_e(\mathfrak{g}) \cap \mathrm{T}_{\mathrm{Q}} = \mathrm{Im}(\mathrm{T}_{\mathrm{P}})\)，且 \(f(\mathrm{G}) = \mathrm{Aut}_e(\mathfrak{g})\)。典则映射

\[
\iota : \mathrm{T} _ {\mathrm{Q}} / \operatorname{Im} (\mathrm{T} _ {\mathrm{P}}) \to \operatorname{Aut} _ {0} (\mathfrak {g}) / \operatorname{Aut} _ {e} (\mathfrak {g})
\]

因此是同构。

\begin{enumerate}
\def\labelenumi{\alph{enumi})}
\setcounter{enumi}{4}
\tightlist
\item
  \(v: \mathrm{G} \to \mathrm{Aut}_e(\mathfrak{g})\) 的核等于 \(\mathrm{T}_{\mathrm{P}} \to \mathrm{T}_{\mathrm{Q}}\) 的核；它同构于
\end{enumerate}

\[
\operatorname{Hom} (\mathrm{P} / \mathrm{Q}, k ^ {*});
\]

这是包含于 G 中心的有限阿贝尔群，其阶整除 \((P:Q)\)；若 k 是代数闭的，则它同构于 P/Q 的对偶（《代数》第VII章§4第8节）。

\(f)\) 令 \(\alpha \in \mathrm{R}\)，\(X_{\alpha} \in \mathfrak{g}^{\alpha}\) 和 \(X_{-\alpha} \in \mathfrak{g}^{-\alpha}\) 满足 \([X_{\alpha}, X_{-\alpha}] = -H_{\alpha}\)，并令 \(\rho_{\alpha}\) 为 \(\mathfrak{sl}(2, k)\) 上相应的 \(\mathfrak{g}\) 表示。若 E 是 \(\mathfrak{g}\)-模，则由§1第4节得到 \(\mathbf{SL}(2, k)\) 在 E 上的表示，因而（由习题25 b)、c)）得到同态

\[
\varphi_ {\alpha}: \mathbf {SL} (2, k) \to \mathrm{G}.
\]

  证明 \(\operatorname{Im}(\varphi_{\alpha})\) 包含 \(T_{P}\) 中形如 \(\lambda \mapsto t^{\lambda(H_{\alpha})}, t \in k^{*}\) 的元素。推出各个 \(\operatorname{Im}(\varphi_{\alpha}), \alpha \in B\) 生成 G（并因而生成 \(T_{P}\)）。特别地，G 由 \(e^{n}\) 生成，其中 \(n \in g^{\alpha}, \alpha \in B \cup -B\)。

\(g\) ) 若 \(\mathbf{G}'\) 是 \(\mathbf{G}\) 的一个子群，满足 \(v(\mathrm{G}') = \mathrm{Aut}_e(\mathfrak{g})\)，则 \(\mathrm{G}' = \mathrm{G}\)（使用 \(f\)）。\(h\) ) 令 \(\mathrm{E}\) 为忠实的 \(\mathfrak{g}\)-模，令 \(\Gamma\) 为 \(\mathrm{P}\) 的权生成的 \(\mathrm{E}\) 的子群。则 \(\mathrm{P} \supset \Gamma \supset \mathrm{Q}\)，参见~习题5。证明典则同态 \(\mathrm{G} \to \mathbf{GL}(\mathrm{E})\) 的核等于 \(T_{P}\) 的子群，该子群由限制在 \(\varphi\) 上为平凡的元素 \(\varGamma\) 组成。特别地，若 \(\varGamma = P\)，同态 \(\mathrm{G} \to \mathbf{GL}(\mathrm{E})\) 是单射。若 \(\varGamma = Q\)，该同态分解为 \(\mathrm{G} \stackrel{v}{\longrightarrow} \operatorname{Aut}_{e}(\mathfrak{g}) \longrightarrow \mathbf{GL}(\mathrm{E})\)，且同态 \(\operatorname{Aut}_{e}(\mathfrak{g}) \to \mathbf{GL}(\mathrm{E})\) 是单射。

\begin{enumerate}
\def\labelenumi{\arabic{enumi})}
\setcounter{enumi}{26}
\tightlist
\item
  令 \(\Omega = \mathrm{P} / \mathrm{Q}\)。若 \(\omega \in \Omega\)，且 \(\mathrm{E}\) 是 \(\mathfrak{g}\)-模，则以 \(\mathrm{E}_{\omega}\) 表示 \(\mathrm{E}^{\lambda}\) 时的 \(\lambda \in \omega\) 的直和。有 \(\mathrm{E} = \bigoplus_{\omega \in \Omega} \mathrm{E}_{\omega}\)。
\end{enumerate}

\begin{enumerate}
\def\labelenumi{\alph{enumi})}
\tightlist
\item
  证明 \(E_{\omega}\) 是 E 的 g-子模。有 \((\mathrm{E}^{*})_{\omega} = (\mathrm{E}_{-\omega})^{*}\) 且
\end{enumerate}

\[
(\mathrm{E} \otimes \mathrm{F}) _ {\omega} = \bigoplus_ {\alpha + \beta = \omega} \mathrm{E} _ {\alpha} \otimes \mathrm{F} _ {\beta}
\]

若 F 是另一个 g-模。

\begin{enumerate}
\def\labelenumi{\alph{enumi})}
\setcounter{enumi}{1}
\item
  令 \(\chi\in\mathrm{Hom}(\varOmega,k^{*})=\mathrm{Ker}(\mathrm{T}_{\mathrm{P}}\to\mathrm{T}_{\mathrm{Q}})\)。将 \(\chi\) 与 \(f:\mathrm{G}\to\mathrm{Aut}_{e}(\mathfrak{g})\) 的核中的一个元素相识别，参见~习题26 e)。证明 \(\chi\) 在 \(E_{\omega}\) 上的作用是比值为 \(\chi(\omega)\) 的伸缩。
\item
  当 \(\mathrm{E}_{\omega}\) 时，\(\mathfrak{g} = \mathfrak{sl}(2,k)\) 是什么？
\end{enumerate}

\exercisesection{\S~8}

\begin{enumerate}
\def\labelenumi{\arabic{enumi})}
\item
  令 \(f\) 为 \(\mathfrak{g}\) 上的不变多项式函数。证明 \(f\) 在 \(\operatorname{Aut}(\mathfrak{g})\) 下不变，当且仅当 \(f|\mathfrak{h}\) 在 \(\operatorname{Aut}(\mathrm{R})\) 下不变。推出，若 \(\mathrm{R}\) 的邓金图有非平凡自同构，则 \(\mathfrak{g}\) 上存在一个在 \(\operatorname{Aut}(\mathfrak{g})\) 下不变的不变多项式函数。
\item
  取 \(\mathfrak{g} = \mathfrak{sl}(3,k)\)。证明 \(x\mapsto \det (x)\) 是 \(\mathfrak{g}\) 上的不变多项式函数，但在 \(\operatorname {Aut}(\mathfrak{g})\) 下不变（使用自同构 \(x\mapsto -^t x\)）。
\item
  令 \(\mathfrak{a}\) 为半单李代数，且 \(s\in \mathrm{Aut}(\mathfrak{a})\)。证明以下条件等价：
\end{enumerate}

\begin{enumerate}
\def\labelenumi{(\roman{enumi})}
\item
  \(s \in \text{Aut}_{0}(\mathfrak{a})\)。
\item
  \(s\) 在 \(\mathrm{U}(\mathfrak{a})\) 的中心上平凡作用。
\item
  对所有 \(x \in \mathfrak{a}\)，存在 \(t \in \mathrm{Aut}_0(\mathfrak{a})\)，使得 \(tx = sx\)。
\end{enumerate}

（使用命题6证明 (iii) \(\Longrightarrow\) (i)。）

\begin{enumerate}
\def\labelenumi{\arabic{enumi})}
\setcounter{enumi}{3}
\item
  证明在命题2的推论2和定理1 (ii) 中，可以只考虑权为根权的表示 \(\rho\)（注意，当 \(k[\mathrm{P}]^{\mathrm{W}}\) 被 \(k[\mathrm{Q}]^{\mathrm{W}}\) 替代时，命题1仍然成立，其中 \(\mathrm{Q}\) 是根权群）。
\item
  沿用§6和§7的记号。令 \(\lambda \in \mathfrak{h}^*\)。若对任意 \(w \in \mathrm{W}\) 且 \(w \neq 1\)，都有 \((\lambda + \rho) - w(\lambda + \rho) \notin \mathrm{Q}_+\)，则 \(\mathrm{Z}(\lambda)\) 是单的。
\end{enumerate}

（使用定理2的推论1 (ii)。）

\begin{enumerate}
\def\labelenumi{\arabic{enumi})}
\setcounter{enumi}{5}
\item
  令 \(\mathfrak{a}\) 为李代数，\(f\) 为 \(\mathfrak{a}\) 上的多项式函数，\(x, y\) 为 \(\mathfrak{a}\) 的两个元素。置 \(f_y = \theta^*(y)f\)（参见~第3节），并以 \(\mathrm{D}_x f\) 表示 \(f\) 在 \(x\) 处的切线性映射（第VII章附录I第2节）。证明 \(f_y(x) = (\mathrm{D}_xf)([x,y])\)。推出，当 \(\mathrm{D}_xf\) 不变时，\(\operatorname{Im}(\operatorname{ad}x)\) 在 \(f\) 上为零。
\item
  *令 \(d_1, \ldots, d_l\) 为代数 I(g) 的特征次数，参见~第V章§5第1节。对任意整数 \(n \geq 0\)，以 \(r_n\) 表示 I(g) 上的 S(g) 的齐次基中次数为 \(n\) 的元素个数，并置 \(r(\mathrm{T}) = \sum_{n=0}^{\infty} r_n \mathrm{T}^n\)。证明
\end{enumerate}

\[
r (\mathrm{T}) = (1 - \mathrm{T}) ^ {- \mathrm{N}} \prod_ {i = 1} ^ {l} (1 - \mathrm{T} ^ {d _ {i}}), \quad \text { where } \mathrm{N} = \dim (\mathfrak {g}).
\]

\begin{enumerate}
\def\labelenumi{\arabic{enumi})}
\setcounter{enumi}{7}
\tightlist
\item
  置 \(l = \operatorname{rk}(\mathfrak{g}), N = \dim(\mathfrak{g})\)。若 \(x \in \mathfrak{g}\)，按下式定义 \(a_i(x)\)，其中 \(0 \leq i \leq N\)：
\end{enumerate}

\[
\det (\mathrm{T} + \operatorname{ad} x) = \sum_ {i = 0} ^ {\mathrm{N}} \mathrm{T} ^ {\mathrm{N} - i} a _ {i} (x).
\]

按上述方式定义的函数 \(a_{i}\) 是次数为 i 的齐次多项式，并且在 \(\operatorname{Aut}(\mathfrak{g})\) 下不变。若 \(x \in h\) 且 \(i \leq N - l\)，则 \(a_{i}(x)\) 是 \(\alpha(x)\)（\(\alpha \in R\)）的第 i 个初等对称函数；特别地，\(a_{N-l}(x) = \prod_{\alpha \in \mathbb{R}} \alpha(x)\)。构造一个例子，使得这些 \(a_{i}\) 不生成在 \(\operatorname{Aut}(\mathfrak{g})\) 下不变的 g 上的多项式函数代数。

\begin{enumerate}
\def\labelenumi{\arabic{enumi})}
\setcounter{enumi}{8}
\tightlist
\item
  记号同第5节。令 \(\lambda \in \mathfrak{h}^*\)，\(z \in \mathbf{Z}\)，\(z'\) 为 z 在 \(z\)\(\mathrm{U}(\mathfrak{g})\) 的主反自同构下的像，令 \(w_0\) 为将 B 变换为 \(-\mathrm{B}\) 的 W 中元素。证明 \(\chi_{\lambda}(z) = \chi_{-w_0\lambda}(z')\)。
\end{enumerate}

（只需对 \(\lambda \in \mathrm{P}_{++}\) 证明此事。考虑 z 在 \(z\)\(\mathrm{E}(\lambda)\) 和 \(\mathrm{E}(\lambda)^*\) 上的作用；使用§7的命题11。）

\begin{enumerate}
\def\labelenumi{\arabic{enumi})}
\setcounter{enumi}{9}
\tightlist
\item
  （本题及以下三题沿用§6的记号。）令 \(\lambda \in \mathfrak{h}^*\)。
\end{enumerate}

\begin{enumerate}
\def\labelenumi{\alph{enumi})}
\tightlist
\item
  令 \(\mathrm{N}, \mathrm{N}'\) 为 \(\mathfrak{g}\) 的 \(\mathrm{Z}(\lambda)\)-子模，满足 \(\mathrm{N}' \subset \mathrm{N}\)，且 \(\mathrm{N} / \mathrm{N}'\) 是单的。证明存在 \(\mu \in \lambda - \mathrm{Q}_+\)，使得 \(\mathrm{N} / \mathrm{N}'\) 同构于 \(\mathrm{E}(\mu)\)（应用§6的定理1），并且 \(\mu + \rho \in \mathrm{W}.(\lambda + \rho)\)（应用定理2的推论1）。
\item
  证明 \(\mathrm{Z}(\lambda)\) 有一个 Jordan-Hölder 序列。（应用 a)，以及 \(\mathrm{Z}(\lambda)\) 的权具有有限重数这一事实。）
\end{enumerate}

\begin{enumerate}
\def\labelenumi{\arabic{enumi})}
\setcounter{enumi}{10}
\tightlist
\item
  令 \(\lambda \in \mathfrak{h}^*\)，令 \(V\) 为 \(\mathfrak{g}\) 的子模且包含于 \(Z(\lambda)\)。证明存在 \(\mu \in \mathfrak{h}^*\)，使得 \(V\) 包含一个简单的 \(\mathfrak{g}\)-子模，它同构于 \(Z(\mu)\)。
\end{enumerate}

（令 A 为满足 V 包含一个同构于 \(\nu \in \mathfrak{h}^*\) 的 \(\mathfrak{g}\)-子模的 \(Z(\nu)\) 的集合。使用第6节的命题6，先证明 \(A \neq \varnothing\)。然后证明 A 是有限的，并取 A 中满足 \(\mu\) 的元素 \((\mu - Q_{+}) \cap A = \{\mu\}\)。）

¶ 12) a) 令 \(\lambda, \mu \in \mathfrak{h}^*\)。从 \(\mathfrak{g}\) 到 \(Z(\mu)\) 的每个非零 \(Z(\lambda)\)-同态都是单射。（使用第6节的命题6。）

\begin{enumerate}
\def\labelenumi{\alph{enumi})}
\setcounter{enumi}{1}
\item
  令 \(r \in \mathbf{N}\)，令 A 为 \(\mathbf{N}^r\) 的有限子集，且 \(m = \mathrm{Card}(\mathrm{A})\)。对于 \(\xi \in \mathbf{N}^r\)，令 \(s(\xi)\) 表示 \(\xi\) 的坐标之和，并令 \(\mathfrak{P}_{\mathrm{A}}(\xi)\) 表示满足 \((n_{\alpha})_{\alpha \in \mathrm{A}}\) 中每个整数 \(\geq 0\) 且 \(\xi = \sum_{\alpha \in \mathrm{A}} n_{\alpha}\alpha\) 的这类族的数目。则有 \(\mathfrak{P}_{\mathrm{A}}(\xi) \leq (s(\xi) + 1)^m\)，对所有 \(\xi \in \mathbf{N}^r\) 成立。（对 \(m\) 作归纳论证。）
\item
  令 \(\lambda, \mu \in \mathfrak{h}^*\)。证明 \(\dim \operatorname{Hom}_{\mathfrak{g}}(\mathrm{Z}(\mu), \mathrm{Z}(\lambda)) \leq 1\)。（令 \(\varphi_1\) 和 \(\varphi_2\) 为非零 \(\mathfrak{g}\)-同态，从 \(\mathrm{Z}(\mu)\) 到 \(\mathrm{Z}(\lambda)\)。若 \(\operatorname{Im}(\varphi_1) = \operatorname{Im}(\varphi_2)\)，则由§6的命题5，\(\varphi_1\) 和 \(\varphi_2\) 线性相关。设 \(\operatorname{Im}(\varphi_1) \neq \operatorname{Im}(\varphi_2)\)。若 \(\mathrm{Z}(\mu)\) 是单的，则 \(\operatorname{Im}(\varphi_1) + \operatorname{Im}(\varphi_2)\) 是直和；推出 \(\mathfrak{P}(\xi + \lambda - \mu) \geq 2\mathfrak{P}(\xi)\) 对所有 \(\xi \in \mathfrak{h}^*\) 都成立，这与 \(b\) 矛盾。在一般情形下，使用相关习题。）
\end{enumerate}

当 \(\dim\operatorname{Hom}_{\mathfrak{g}}(\mathrm{Z}(\mu),\mathrm{Z}(\lambda))=1\) 时，滥用记号写作 \(\mathrm{Z}(\mu)\subset\mathrm{Z}(\lambda)\)。d) 令 \(\nu\in\mathfrak{h}^{*}\)。使得 \(\lambda\in\mathfrak{h}^{*}\) 的 \(\mathrm{Z}(\lambda-\nu)\subset\mathrm{Z}(\lambda)\) 的集合在 \(\mathfrak{h}^{*}\) 的 Zariski 拓扑中是闭集。

¶ 13) a) 令 \(\mathfrak{a}\) 为幂零李代数，\(x \in \mathfrak{a}, n \in \mathbf{N}, p \in \mathbf{N}\)。存在 \(l \in \mathbf{N}\)，使得对所有 \(x^{l}y_{1} \ldots y_{n} \in \mathrm{U}(\mathfrak{a})x^{p}\)，都有 \(y_{1}, \ldots, y_{n} \in \mathfrak{a}\)。

\begin{enumerate}
\def\labelenumi{\alph{enumi})}
\setcounter{enumi}{1}
\tightlist
\item
  令 \(\lambda, \mu \in \mathfrak{h}^*\)，且 \(\alpha \in \mathrm{B}\) 满足
\end{enumerate}

\[
\mathrm{Z} (s _ {\alpha} \mu - \rho) \subset \mathrm{Z} (\mu - \rho) \subset \mathrm{Z} (\lambda - \rho).
\]

假设 \(\lambda \in \mathrm{P}\)。令 \(p = \lambda (H_{\alpha})\in \mathbf{Z}\)。证明：

\(b_{1})\) 若 \(p \leq 0\)，则 \(\mathrm{Z}(\lambda - \rho) \subset \mathrm{Z}(s_{\alpha}\lambda - \rho)\)。

\[
b _ {2}) \text {   If   } p > 0, \text {   then   } \mathrm{Z} (s _ {\alpha} \mu - \rho) \subset \mathrm{Z} (s _ {\alpha} \lambda - \rho) \subset \mathrm{Z} (\lambda - \rho).
\]

（使用 \(a\)，以及§6命题6的推论1。）

\begin{enumerate}
\def\labelenumi{\alph{enumi})}
\setcounter{enumi}{2}
\tightlist
\item
  令 \(\lambda \in \mathfrak{h}^*\)，\(\alpha \in \mathrm{R}\)，且 \(m = \lambda(H_{\alpha})\)。假设 \(m \in \mathbf{N}\)。证明
\end{enumerate}

\[
\mathrm{Z} (s _ {\alpha} \lambda - \rho) \subset \mathrm{Z} (\lambda - \rho).
\]

(先对 \(\lambda \in \mathrm{P}\) 使用 \(b\) 证明此事，然后在一般情形下使用习题12 \(d\)。) \(^{17}\)

\begin{enumerate}
\def\labelenumi{\arabic{enumi})}
\setcounter{enumi}{13}
\item
  *令 \(\mathfrak{a}\) 为半单李代数，令 \(Z(\mathfrak{a})\) 为 \(U(\mathfrak{a})\) 的中心。证明 \(U(\mathfrak{a})\) 是自由的 \(Z(\mathfrak{a})\)-模。（注意 \(grU(\mathfrak{a})\) 同构于 \(S(\mathfrak{a})\) 和 \(S(\mathfrak{a}^*)\)，并使用第3节的注2。）*
\item
  令 \(x\) 为 \(\mathfrak{g}\) 的一个可对角化元素（§3，习题10），令 \(y\) 为 \(\mathfrak{g}\) 的一个半单元素，满足对 \(f(x) = f(y)\) 上的每个不变多项式函数 \(f\) 都有 \(\mathfrak{g}\)。证明存在 \(s \in \mathrm{Aut}_e(\mathfrak{g})\)，使得 \(sy = x\)。（注意，ad \(x\) 和 ad \(y\) 有相同的特征多项式，参见~习题8，因而 y 是可对角化的。接着使用§3的习题10，将情形化归为 x 和 y 包含于 h 中，再使用定理1 (i) 和引理6证明 x 和 y 在 W 下共轭。）
\item
  令 \(\mathfrak{a}\) 为半单李代数，\(x\) 为 \(\mathfrak{a}\) 的一个元素，\(x_{s}\) 为 \(x\) 的半单部分。证明，若 \(f\) 是 \(\mathfrak{a}\) 上的不变多项式函数，则 \(f(x) = f(x_{s})\)。（将情形化归为 \(f\) 形如 \(x \mapsto \operatorname{Tr} \rho(x)^n\) 的情形。）
\item
  假设 \(k\) 是代数闭域，并置 \(\mathrm{G} = \operatorname{Aut}_e(\mathfrak{g})\)。令 \(x, y \in \mathfrak{g}\)。证明以下条件等价：
\end{enumerate}

\begin{enumerate}
\def\labelenumi{(\roman{enumi})}
\item
  \(x\) 和 \(y\) 的半单部分是 G-共轭的。
\item
  对 \(f\) 上的每个不变多项式函数 \(\mathfrak{g}\)，都有 \(f(x) = f(y)\)。（使用习题15和16。）
\end{enumerate}

\begin{enumerate}
\def\labelenumi{\arabic{enumi})}
\setcounter{enumi}{17}
\tightlist
\item
  令 \(\mathfrak{a}\) 为半单李代数，\(l = \mathrm{rk}(\mathfrak{a})\)，I 为 \(\mathfrak{a}\) 上不变多项式函数的代数，且 \(\mathrm{P}_1, \ldots, \mathrm{P}_l\) 为 I 的齐次元素，并且作为代数生成 I。\(\mathrm{P}_i\) 定义多项式映射 \(\mathrm{P} : \mathfrak{a} \to k^l\)。若 \(x \in \mathfrak{a}\)，以 \(\mathrm{D}_x\mathrm{P} : \mathfrak{a} \to k^l\) 表示 P 在 \(x\) 处的切线性映射（第VII章附录I第2节）。
\end{enumerate}

\begin{enumerate}
\def\labelenumi{\alph{enumi})}
\tightlist
\item
  令 \(\mathfrak{h}\) 为 \(\mathfrak{a}\) 的 Cartan 子代数，令 \(x \in \mathfrak{h}\)。证明以下条件等价：
\end{enumerate}

\begin{enumerate}
\def\labelenumi{(\roman{enumi})}
\item
  \(D_{x}P|\mathfrak{h}\) 是从 \(\mathfrak{h}\) 到 \(k^{l}\) 的同构；
\item
  \(x\) 是正则的。
\end{enumerate}

（将情形化归为分裂情形。取 \(\mathfrak{h}\) 的一个基，并记 \(d(x)\) 关于此基的矩阵的行列式为 \(\mathrm{D}_x\mathrm{P}|\mathfrak{h}\)。利用第V章§5第4节的命题5证明，存在 \(c \in k^*\)，使得 \(d(x)^2 = c \prod_{\alpha \in \mathrm{R}} \alpha(x)\)，

其中 \(\alpha\) 属于 \(R\) 的根集 \((\mathfrak{a},\mathfrak{h})\)。）

若满足这些条件，则 \(\mathfrak{a} = \mathfrak{h} \oplus \operatorname{Im} \operatorname{ad}(x)\)，且 \(\operatorname{Ker} D_x P = \operatorname{Im} \operatorname{ad}(x)\)（使用习题6证明 \(D_x P\) 在 \(\operatorname{Im} \operatorname{ad}(x)\) 上为零）。

\begin{enumerate}
\def\labelenumi{\alph{enumi})}
\setcounter{enumi}{1}
\tightlist
\item
  证明满足 \(x \in a\) 秩为 l 的 \(D_{x}P\) 的集合是 a 在 Zariski 拓扑中的稠密开子集。
\end{enumerate}

\exercisesection{\S~9}

以下所考虑的所有 g-模都假设为有限维的。

\begin{enumerate}
\def\labelenumi{\arabic{enumi})}
\item
  若 \(m\) 是满足 \(\geq 0\) 的整数，则有 \(\dim \mathrm{E}(m\rho) = (m + 1)^{\mathrm{N}}\)，其中 \(\mathrm{N} = \operatorname{Card}(\mathrm{R}_{+})\)。
\item
  证明 \(d\) 上存在唯一的多项式函数 \(\mathfrak{h}^*\)，使得对所有 \(d(\lambda) = \dim \mathrm{E}(\lambda)\)，都有 \(\lambda \in \mathrm{P}_{++}\)；其次数为 \(\operatorname{Card}(\mathrm{R}_+)\)。有
\end{enumerate}

\[
d (w \lambda - \rho) = \varepsilon (w) d (\lambda - \rho) \quad \text { if } w \in W, \lambda \in \mathfrak {h} ^ {*}.
\]

特别地，函数 \(\lambda \mapsto d(\lambda - \rho)^2\) 在 W 下不变。推出 \(u\) 的中心中存在唯一元素 \(\mathrm{U}(\mathfrak{g})\)，使得对所有 \(\chi_{\lambda}(u) = d(\lambda)^2\)，都有 \(\lambda \in \mathfrak{h}^*\)（应用§8第5节的定理2）。当 \(\mathfrak{g} = \mathfrak{sl}(2,k)\) 时，有 \(u = C + 1\)，其中 \(C\) 是§1习题1中定义的元素。

\begin{enumerate}
\def\labelenumi{\arabic{enumi})}
\setcounter{enumi}{2}
\tightlist
\item
  令 \(k_{1}, \ldots, k_{l}\) 为 W 的不变量代数的特征次数（参见~第V章§5）。
\end{enumerate}

\begin{enumerate}
\def\labelenumi{\alph{enumi})}
\item
  证明对所有 \(j \geq 1\)，满足 \(k_{i} > j\) 的 i 的个数等于满足 \(\alpha \in R_{+}\) 的 \(\langle \rho, H_{\alpha} \rangle = j\) 的个数。（将情形化归为 R 不可约的情形，并使用第VI章§4习题6 c)。）（参见§5习题5 g)。
\item
  推出公式
\end{enumerate}

\[
\prod_ {\alpha \in \mathrm{R} _ {+}} \langle \rho , H _ {\alpha} \rangle = \prod_ {i = 1} ^ {l} (k _ {i} - 1)!.
\]

¶4) 假设 g 是单的，记 \(\gamma\) 为 \(R_{+}\) 中满足 \(H_{\gamma}\) 是 \(R^{\vee}\) 的最高根的元素；写作 \(H_{\gamma} = \sum_{\alpha \in B} n_{\alpha} H_{\alpha}\)。有 \(\langle \rho, H_{\gamma} \rangle = \sum n_{\alpha} = h - 1\)，其中 h 是 R 的 Coxeter 数（第VI章§1第11节命题31）。

\begin{enumerate}
\def\labelenumi{\alph{enumi})}
\tightlist
\item
  令 \(\alpha \in \mathrm{B}\)。证明对所有 \(\beta \in \mathrm{R}_{+}\)，
\end{enumerate}

\[
\langle \varpi_ {\alpha} + \rho , H _ {\beta} \rangle \leq h + n _ {\alpha} - 1,
\]

并且当 \(\beta = \gamma\) 时取等号。推出 \(\dim \mathrm{E}(\varpi_{\alpha})\) 的每个素因子都 \(\leq h + n_{\alpha} - 1\)。

\begin{enumerate}
\def\labelenumi{\alph{enumi})}
\setcounter{enumi}{1}
\tightlist
\item
  假设 \(\varpi_{\alpha}\) 不是极小权，即~\(n_{\alpha} \geq 2\)。令 \(m \in [2, n_{\alpha}]\) 且 \(p = h + m - 1\)。验证（参见~第VI章图版），存在 \(\beta \in \mathrm{R}_{+}\)，使得 \(\langle \varpi_{\alpha}, H_{\beta} \rangle = n_{\alpha}\)，且 \(\langle \rho, H_{\beta} \rangle = h - 1 - (n_{\alpha} - m)\)，从而 \(\langle \varpi_{\alpha} + \rho, H_{\beta} \rangle = p\)。推出，若 \(p\) 是素数，则 \(p\) 整除 \(\dim \mathrm{E}(\varpi_{\alpha})\)。（注意，对于 \(p\)，\(\langle \rho, H_{\beta} \rangle\) 不整除任何 \(\beta \in \mathrm{R}_{+}\)，参见~习题3。）
\end{enumerate}

\(c\) ) 当 \(\mathfrak{g}\) 为 \(\mathrm{G}_2\) 型（分别为 \(\mathrm{F}_4,\mathrm{E}_8\) 型）时，有 \(h = 6\)（分别为 12、30），且 \(\dim \mathrm{E}(\varpi_{\alpha})\) 被 7 整除（分别被 13、31 整除）；当 \(\mathfrak{g}\) 为 \(\mathrm{E}_6\) 型（分别为 \(\mathrm{E}_7\) 型），且 \(\varpi_{\alpha}\) 不是极小权时，\(\dim \mathrm{E}(\varpi_{\alpha})\) 被 13 整除（分别被 19 整除）。

¶ 5) a) 令 \(\alpha \in \mathrm{R}, x \in \mathfrak{g}^{\alpha}, y \in \mathfrak{g}^{-\alpha}\)，且 E 为 \(\mathfrak{g}\)-模。证明对所有 \(\lambda \in \mathrm{P}\)，

\[
\operatorname{Tr} ((x y) _ {\mathrm{E}} | \mathrm{E} ^ {\lambda}) = \operatorname{Tr} ((x y) _ {\mathrm{E}} | \mathrm{E} ^ {\lambda + \alpha}) + \lambda ([ x, y ]) \dim \mathrm{E} ^ {\lambda}.
\]

推出

\[
\sum_ {\lambda \in P} \lambda ([ x, y ]) \dim E ^ {\lambda}. e ^ {\lambda} = (1 - e ^ {- \alpha}) \sum_ {\lambda \in P} \operatorname{Tr} ((x y) _ {E} | E ^ {\lambda}). e ^ {\lambda}.
\]

\begin{enumerate}
\def\labelenumi{\alph{enumi})}
\setcounter{enumi}{1}
\tightlist
\item
  在 \(\mathfrak{h}^*\) 上取一个非退化的 W-不变对称双线性形式 \(\langle ,\rangle\)。令 \(\Delta\) 为向量空间 \(k[\mathrm{P}]\) 的自同态，使得对所有 \(\Delta(e^{\mu}) = \langle \mu ,\mu \rangle e^{\mu}\)，都有 \(\mu \in \mathrm{P}\)；若 \(a,b\in k[\mathrm{P}]\)，置
\end{enumerate}

\[
\Delta^ {\prime} (a, b) = \Delta (a b) - a \Delta (b) - b \Delta (a).
\]

证明

\[
\Delta (\mathrm{J} (e ^ {\mu})) = \langle \mu , \mu \rangle \mathrm{J} (e ^ {\mu}) \quad \text { for } \mu \in \mathrm{P},
\]

\[
\Delta^ {\prime} (e ^ {\lambda}, e ^ {\mu}) = 2 \langle \lambda , \mu \rangle e ^ {\lambda + \mu} \quad \mathrm{for} \lambda , \mu \in \mathrm{P},
\]

\[
\Delta^ {\prime} (a b, c) = a \Delta^ {\prime} (b, c) + b \Delta^ {\prime} (a, c) \quad \text { for } a, b, c \in k [ P ].
\]

\(c)\) 令 \(\lambda \in \mathrm{P}_{++}\)，\(c_{\lambda} = \operatorname{ch}(\operatorname{E}(\lambda))\) 且 \(d = \mathrm{J}(e^{\rho})\)。证明

\[
\Delta (c _ {\lambda} d) = \langle \lambda + \rho , \lambda + \rho \rangle c _ {\lambda} d.
\]

（使用 \(a\)、\(b\)、§6命题7的推论，以及第VI章§3第3节的公式 (3)。）

\(d)\) 由上述结果和§7第2节命题5 (iii) 推出 H. Weyl 公式的另一种证明。

\begin{enumerate}
\def\labelenumi{\alph{enumi})}
\setcounter{enumi}{4}
\tightlist
\item
  对所有 \(\lambda \in \mathfrak{h}^*\)，置 \(\dim \mathrm{E}^{\lambda} = m(\lambda)\)。由 \(a\) 推出
\end{enumerate}

\[
\begin{array}{c} \operatorname{Tr} ((x y) _ {\mathrm{E}} | \mathrm{E} ^ {\lambda}) = \sum_ {i = 0} ^ {+ \infty} (\lambda + i \alpha) ([ x, y ]) m (\lambda + i \alpha) \\ \sum_ {i = - \infty} ^ {+ \infty} (\lambda + i \alpha) ([ x, y ]) m (\lambda + i \alpha) = 0. \end{array}
\]

\(f)\) 令 \(\langle\cdot,\cdot\rangle\) 为 \(\mathfrak{g}\) 上的非退化不变对称双线性形式，其在 \(\mathfrak{h}\) 上的限制是上面所选形式的逆。令 \(\Gamma\) 为相应的 Casimir 元。假设 E 是单的；置 \(\Gamma_{\mathrm{E}}=\gamma.1\)，其中 \(\gamma\in k\)。利用 \(e\) 和§2第3节命题6，证明

\[
\gamma m (\lambda) = \langle \lambda , \lambda \rangle m (\lambda) + \sum_ {\alpha \in \mathrm{R}} \sum_ {i = 0} ^ {+ \infty} \langle \lambda + i \alpha , \alpha \rangle m (\lambda + i \alpha)
\]

对所有 \(\lambda \in \mathfrak{h}^*\)，并进而证明

\[
\begin{array}{l} \gamma m (\lambda) = \langle \lambda , \lambda \rangle m (\lambda) + \sum_ {\alpha \in \mathrm{R} _ {+}} m (\lambda) \langle \lambda , \alpha \rangle + 2 \sum_ {\alpha \in \mathrm{R} _ {+}} \sum_ {i = 1} ^ {+ \infty} m (\lambda + i \alpha) \langle \lambda + i \alpha , \alpha \rangle \\ = \langle \lambda , \lambda + 2 \rho \rangle m (\lambda) + 2 \sum_ {\alpha \in \mathrm{R} _ {+}} \sum_ {i = 1} ^ {+ \infty} m (\lambda + i \alpha) \langle \lambda + i \alpha , \alpha \rangle . \end{array}
\]

\(g)\) 继续假设 E 是单的；令 \(\omega\) 为其最高权。由 \(f)\) 推出，对所有 \(\lambda \in \mathfrak{h}^*\)，

\[
(\langle \omega + \rho , \omega + \rho \rangle - \langle \lambda + \rho , \lambda + \rho \rangle) m (\lambda) = 2 \sum_ {\alpha \in \mathrm{R} _ {+}} \sum_ {i = 1} ^ {+ \infty} m (\lambda + i \alpha) \langle \lambda + i \alpha , \alpha \rangle .
\]

（回忆由§7的命题5，若 \(\langle\omega+\rho,\omega+\rho\rangle>\langle\lambda+\rho,\lambda+\rho\rangle\) 是 E 的不同于 \(\lambda\) 的权，则 \(\omega\)。因此，上式给出了逐步计算 \(m(\lambda)\) 的方法。） \(^{18}\)

\begin{enumerate}
\def\labelenumi{\arabic{enumi})}
\setcounter{enumi}{5}
\tightlist
\item
  令 \(x \mapsto x^*\) 为 \(k[\mathrm{P}]\) 上的对合，将所有 \(e^p\) 的 \(e^{-p}\) 变为 \(p \in \mathrm{P}\)。\(a)\) 置 \(\mathrm{D} = d^* d = \prod_{\alpha \in \mathrm{R}} (1 - e^\alpha)\)。证明 \(d^* = (-1)^{\mathrm{N}} d\)，其中 \(\mathrm{N} = \mathrm{Card}(\mathrm{R}_+)\)，从而 \(\mathrm{D} = (-1)^{\mathrm{N}} d^2\)。
\end{enumerate}

\begin{enumerate}
\def\labelenumi{\alph{enumi})}
\setcounter{enumi}{1}
\tightlist
\item
  按下列公式在 \(\varepsilon\) 上定义两个线性形式 \(k[\mathrm{P}]\) 和 I：
\end{enumerate}

\[
\varepsilon (1) = 1, \quad \varepsilon (e ^ {p}) = 0 \quad \text { if } p \in \mathrm{P} - \{0 \}
\]

以及

\[
\mathrm{I} (f) = \frac {1}{m} \varepsilon (\mathrm{D}. f), \quad \mathrm{where} m = \mathrm{Card} (\mathrm{W}).
\]

利用公式 \(d = \mathrm{J}(e^{\rho})\)，证明 \(\mathrm{I}(1) = 1\)。

\begin{enumerate}
\def\labelenumi{\alph{enumi})}
\setcounter{enumi}{2}
\tightlist
\item
  令 \(\lambda \in \mathrm{P}_{++}\)，且 \(c_{\lambda} = \operatorname{ch} \mathrm{E}(\lambda) = \mathrm{J}(e^{\lambda + \rho}) / d\)。证明若 \(\mathrm{I}(c_{\lambda}) = 0\)，则 \(\lambda \neq 0\)。（方法同 \(b\)。）
\end{enumerate}

\(d\) ) 证明 I 在 \(\mathbf{Z}[\mathrm{P}]^{\mathrm{W}} = \operatorname {ch}\operatorname {R}(\mathfrak{g})\) 的子代数 \(k[\mathrm{P}]\) 上取整数值。若 E 是 \(\mathfrak{g}\)-模，则 E 中 \(\mathfrak{g}\) 的不变子空间的维数等于 I(ch E)。（将情形化归为 E 为单的情形，并使用 \(b\)）和 \(c\)。）\(e\) ) 有 \(\dim \mathrm{E} = \sum_{\lambda \in \mathrm{P}_{++}}\mathrm{I}(c_\lambda^*\mathrm{ch}\mathrm{E})d(\lambda)\)，其中 \(d(\lambda) = \dim \mathrm{E}(\lambda)\)。特别地：

\[
\mathrm{I} (c _ {\lambda} ^ {*} c _ {\mu}) = \delta_ {\lambda \mu} \quad \text { if } \lambda , \mu \in \mathrm{P} _ {+ +}.
\]

\(f)\) 若 \(\lambda, \mu, \nu \in \mathrm{P}_{++}\)，则命题2中的整数 \(m(\lambda, \mu, \nu)\) 等于 \(\mathrm{I}(c_{\lambda} c_{\mu} c_{\nu}^{*})\)。推出恒等式

\[
d (\lambda) d (\mu) = \sum_ {\nu \in \mathrm{P} _ {+ +}} m (\lambda , \mu , \nu) d (\nu) \quad \lambda , \mu , \nu \in \mathrm{P} _ {+ +}.
\]

（将 e) 应用于 g-模 \(\mathrm{E} = \mathrm{E}(\lambda) \otimes \mathrm{E}(\mu).\)

¶ 7) 沿用定理2证明中的记号。

\begin{enumerate}
\def\labelenumi{\alph{enumi})}
\tightlist
\item
  证明
\end{enumerate}

\[
f _ {\rho} (\mathrm{J} (e ^ {\mu})) = \prod_ {\alpha \in \mathrm{R} _ {+}} (e ^ {(\mu | \alpha) \mathrm{T} / 2} - e ^ {- (\mu | \alpha) \mathrm{T} / 2}).
\]

\begin{enumerate}
\def\labelenumi{\alph{enumi})}
\setcounter{enumi}{1}
\tightlist
\item
  取 \((\cdot |\cdot)\) 为典则双线性形式 \(\varPhi_{\mathrm{R}}\)（第VI章§1第12节）。证明
\end{enumerate}

\[
  f _ {\rho} (\mathrm{J} (e ^ {\mu})) \equiv d _ {\mu} \mathrm{T} ^ {\mathrm{N}} \left(1 + \frac {\mathrm{T} ^ {2}}{48} (\mu | \mu)\right) \pmod {\mathrm{T} ^ {\mathrm{N} + 3} \mathbf {R} [ [ \mathrm{T} ] ]},
\]

其中 \(d_{\mu} = \prod_{\alpha \in \mathrm{R}_{+}}(\mu |\alpha)\)。

\begin{enumerate}
\def\labelenumi{\alph{enumi})}
\setcounter{enumi}{2}
\tightlist
\item
  由 b) 以及等式 \(\mathrm{J}(e^{\lambda+\rho})=\mathrm{ch}(\mathrm{E}).\mathrm{J}(e^{\rho})\)，推出公式
\end{enumerate}

\[
  \sum_ {\mu \in \mathrm{P}} (\mu | \rho) ^ {2} \dim \mathrm{E} ^ {\mu} = \frac {\dim \mathrm{E}}{24} (\lambda | \lambda + 2 \rho).
\]

\(d\) ) 假设 \(\mathfrak{g}\) 是单的。证明 \((\rho |\rho) = \dim \mathfrak{g} / 24\)。（在 \(c\) 为 R 的最高根时应用 \(\lambda\)，并使用§6的习题3。）

¶8) 令 \(\psi\) 为 \(\mathfrak{h}^*\) 上次数为 \(r\) 的多项式函数。证明存在唯一的 \(\varPsi\) 上的多项式函数 \(\mathfrak{h}^*\)，它在 W 下不变，次数 \(\leq r\)，并且满足

\[
\sum_ {\mu \in \mathrm{P}} \psi (\mu) \dim \mathrm{E} ^ {\mu} = \varPsi (\lambda + \rho) \dim \mathrm{E}
\]

对任意最高权为 \(\lambda\) 的单 g-模 E。

（先处理 \(\psi(\mu)=(\mu|\nu)^{r}\) 的情形，其中 \(\nu\in P\) 不与任何根正交；为此使用定理2的证明中的同态 \(f_{\nu}\) 以及第V章§5第4节命题5 (i)。）

\begin{enumerate}
\def\labelenumi{\arabic{enumi})}
\setcounter{enumi}{8}
\tightlist
\item
  在 \(\mathfrak{g}\) 为 \(\mathrm{A}_2\) 型代数的情形，沿用第VI章图版I的记号。令 \(n, p\) 为满足 \(\geq 0\) 的整数。
\end{enumerate}

\[
a) \mathfrak {P} (n \alpha_ {1} + p \alpha_ {2}) = 1 + \inf (n, p).
\]

\begin{enumerate}
\def\labelenumi{\alph{enumi})}
\setcounter{enumi}{1}
\tightlist
\item
  令 \(\lambda = n\varpi_1 + p\varpi_2\)。则 \(\dim \mathrm{E}(\lambda) = \frac{1}{2}(n + 1)(p + 1)(n + p + 2)\)。若 \(\mathrm{E}(\lambda)\) 不是根权，即~\(\lambda\)，则 \(n \not\equiv p \pmod{3}\) 的权 0 的重数为 0；若 \(\lambda\) 是根权，则该重数为 \(1 + \inf(n,p)\)。
\end{enumerate}

\begin{enumerate}
\def\labelenumi{\arabic{enumi})}
\setcounter{enumi}{9}
\tightlist
\item
  在代数为 \(B_{2}\) 型的情形，沿用第VI章图版II的记号。令 n, p 为满足 \(\geq 0\) 的整数。
\end{enumerate}

\begin{enumerate}
\def\labelenumi{\alph{enumi})}
\tightlist
\item
  有
\end{enumerate}

\[
\begin{array}{c} \mathfrak {P} (n \alpha_ {1} + p (\alpha_ {1} + 2 \alpha_ {2})) = 1 + \frac {1}{2} p (p + 3) \\ \mathfrak {P} (n \alpha_ {2} + p (\alpha_ {1} + \alpha_ {2})) = [ p ^ {2} / 4 ] + p + 1 \\ \mathfrak {P} (n (\alpha_ {1} + \alpha_ {2}) + p (\alpha_ {1} + 2 \alpha_ {2})) = [ n ^ {2} / 4 ] + n + 1 + n p + \frac {1}{2} p (p + 3). \end{array}
\]

\begin{enumerate}
\def\labelenumi{\alph{enumi})}
\setcounter{enumi}{1}
\tightlist
\item
  令 \(\lambda = n(\alpha_{1} + \alpha_{2}) + p(\alpha_{1} + 2\alpha_{2})\)。\(\operatorname{E}(\lambda)\) 中权 0 的重数为
\end{enumerate}

\[
[ n / 2 ] + 1 + n p + p.
\]

\begin{enumerate}
\def\labelenumi{\arabic{enumi})}
\setcounter{enumi}{10}
\item
  假设 \(\mathfrak{g}\) 不是秩为 1 的代数的乘积。令 \(n\in \mathbf{N}\)。证明存在一个单 \(\mathfrak{g}\)-模，其某个权的重数 \(\geq n\)。（反设不然。令 \(E_{\lambda}\) 为最高权 \(\lambda \in P_{++}\) 的单 g-模。当 \(\dim E_{\lambda}\)\(E_{\lambda}\) 时，比较 \((\lambda|\lambda) \to \infty\) 的维数和其不同权的数目（使用定理2的记号）。）
\item
  令 \(U^{0}\) 为 \(\mathrm{U}(\mathfrak{g})\) 中 h 的交换子。若 g 的秩为 1，则 \(U^{0}\) 是交换的。若 g 的秩 \(\geq 2\)，则 \(U^{0}\) 有任意大的有限维单表示。（使用习题11和§7习题23 a)。）
\item
  令 R 为向量空间 V 中的根系。若 R 是两个根系 \(v_{1}, v_{2}\) 和 \(R_{1}\) 的直和（第VI章§1第2节），且 \(R_{2}\) 属于 Rᵢ 生成的 V 的向量子空间（\(v_{i}\)），则称 V 的两个元素 \(R_{i}, i = 1, 2\) 不相交。证明，属于 R 的同一室的 \(V_{R}\) 中两个元素不相交，当且仅当它们正交。
\end{enumerate}

¶ 14) a) 令 \(\mu, \nu \in \mathrm{P}_{++}\)，令 \(\gamma\) 为 \(\mathrm{E}(\mu)\) 的一个权。令 \(\rho_{\mu}\) 为 \(\mathfrak{g}\) 在 \(\mathrm{E}(\mu)\) 上的表示。令 \(X_{\alpha} \in \mathfrak{g}^{\alpha} - \{0\}, Y_{\alpha} \in \mathfrak{g}^{-\alpha} - \{0\}\)。若 \(\alpha \in \mathrm{B}\)，置 \(\nu_{\alpha} = \nu(H_{\alpha})\)，并

\[
\mathrm{E} ^ {+} (\mu , \gamma , \nu) = \mathrm{E} (\mu) ^ {\gamma} \cap \bigcap_ {\alpha \in \mathrm{B}} \mathrm{Ker} \rho_ {\mu} (X _ {\alpha}) ^ {\nu_ {\alpha} + 1},
\]

\[
\mathrm{E} ^ {-} (\mu , \gamma , \nu) = \mathrm{E} (\mu) ^ {\gamma} \cap \bigcap_ {\alpha \in \mathrm{B}} \operatorname{Ker} \rho_ {\mu} (Y _ {\alpha}) ^ {\nu_ {\alpha} + 1},
\]